
%====================================================================
% MOVEMENT VI: GRAVITATIONAL S-DUALITY
%====================================================================

\subsection{Movement VI: gravitational $S$-duality}
\label{subsec:gravity-s-duality}
\index{S-duality@$S$-duality!gravitational|textbf}
\index{Verdier duality!gravitational $S$-duality}
\index{Koszul involution!gravitational}

The Virasoro Koszul duality $\mathrm{Vir}_c^! = \mathrm{Vir}_{26-c}$
(Proposition~\ref{prop:gravity-koszul-dual}) is not merely a
formal consequence of the bar-cobar adjunction: it is an exact
duality of the gravitational theory.  Movement~VI derives this
duality at the level of the full MC element, the shadow tower,
the genus expansion, and the holographic connection.  The
organizing principle: the Verdier involution
$\mathbb{D}$ on $\operatorname{Ran}(X)$ sends
$\Theta_{\mathrm{Vir}_c}$ to $\Theta_{\mathrm{Vir}_{26-c}}$,
and every gravitational observable transforms predictably under
$c \mapsto 26 - c$.

\subsubsection*{The Verdier intertwining theorem}

\begin{theorem}[Verdier intertwining for Virasoro;
\ClaimStatusProvedHere]
\label{thm:gravity-verdier-intertwining}
\index{Verdier intertwining!Virasoro}
Let $\mathbb{D} \colon D(\operatorname{Ran}(X))
\xrightarrow{\;\sim\;} D(\operatorname{Ran}(X))^{\mathrm{op}}$
denote the Verdier duality functor.  Then:
\begin{enumerate}[label=\textup{(\roman*)}]
\item $\mathbb{D}(\barB(\mathrm{Vir}_c))
  \simeq \barB(\mathrm{Vir}_{26-c})$ as dg coalgebras on
  $\operatorname{Ran}(X)$.
\item $\mathbb{D}(\Theta_{\mathrm{Vir}_c})
  = \Theta_{\mathrm{Vir}_{26-c}}$ as MC elements in the
  modular cyclic deformation complex
  $\mathrm{Def}_{\mathrm{cyc}}^{\mathrm{mod}}(\mathrm{Vir}_{26-c})$,
  where both sides are identified via the level-swap isomorphism
  $\mathbb{D} \colon
  \mathrm{Def}_{\mathrm{cyc}}^{\mathrm{mod}}(\mathrm{Vir}_c)
  \xrightarrow{\;\sim\;}
  \mathrm{Def}_{\mathrm{cyc}}^{\mathrm{mod}}(\mathrm{Vir}_{26-c})$
  induced by $\mathbb{D}(\barB(\mathrm{Vir}_c))
  \simeq \barB(\mathrm{Vir}_{26-c})$.
\item For every $r \ge 2$,
  $\mathbb{D}(\operatorname{Sh}_r(\mathrm{Vir}_c))
  = \operatorname{Sh}_r(\mathrm{Vir}_{26-c})$.
\end{enumerate}
\end{theorem}

\begin{proof}
(i).\ The bar construction $\barB(\cA)$ is an algebra over
$\operatorname{FCom}$ (Theorem~\ref{thm:bar-modular-operad}).
Verdier duality $\mathbb{D}$ sends the chiral pairing
$\langle -, - \rangle_{\mathrm{Vir}_c}$ to the dual pairing
with $c \mapsto 26 - c$.  Explicitly: the Virasoro Killing form
on weight-$n$ states is
$\langle L_{-n}, L_{-n} \rangle = \frac{c}{12}(n^3 - n)$;
under $\mathbb{D}$ this maps to
$\frac{26-c}{12}(n^3 - n)$, which is
$\langle L_{-n}, L_{-n} \rangle_{\mathrm{Vir}_{26-c}}$.
Since the bar differential is determined by the OPE and the
pairing, and $\mathbb{D}$ interchanges them, we obtain
$\mathbb{D}(\barB(\mathrm{Vir}_c))
\simeq \barB(\mathrm{Vir}_{26-c})$.

(ii).\ The MC element $\Theta_\cA := D_\cA - d_0$ is
bar-intrinsic (Theorem~\ref{thm:mc2-bar-intrinsic}).  Part~(i)
gives $\mathbb{D}(D_{\mathrm{Vir}_c})
= D_{\mathrm{Vir}_{26-c}}$ and
$\mathbb{D}(d_0) = d_0$.  Hence
$\mathbb{D}(\Theta_{\mathrm{Vir}_c})
= D_{\mathrm{Vir}_{26-c}} - d_0
= \Theta_{\mathrm{Vir}_{26-c}}$.

(iii).\ The shadows are finite-order projections:
$\operatorname{Sh}_r(\cA)
= \pi_{\le r}(\Theta_\cA)$.  Since $\mathbb{D}$ preserves
arity, $\pi_{\le r} \circ \mathbb{D}
= \mathbb{D} \circ \pi_{\le r}$, and the claim follows from~(ii).
\end{proof}

\subsubsection*{The complementarity constant}

\begin{proposition}[Complementarity at $K = 13$;
\ClaimStatusProvedHere]
\label{prop:gravity-complementarity-constant}
\index{complementarity constant!gravitational}
The complementarity constant of the Virasoro Koszul pair
$(\mathrm{Vir}_c, \mathrm{Vir}_{26-c})$ is
\begin{equation}\label{eq:gravity-complementarity-K}
K
\;=\;
\kappa(\mathrm{Vir}_c) + \kappa(\mathrm{Vir}_{26-c})
\;=\;
\frac{c}{2} + \frac{26-c}{2}
\;=\; 13.
\end{equation}
This is independent of~$c$: the total curvature of the Koszul
pair is a universal constant.
\end{proposition}

\begin{proof}
$\kappa(\mathrm{Vir}_c) = c/2$
(Theorem~\ref{thm:gravity-koszul-triangle}).
The Koszul dual has $\kappa(\mathrm{Vir}_{26-c}) = (26-c)/2$.
Sum: $(c + 26 - c)/2 = 13$.
\end{proof}

The gravitational meaning: the total one-loop anomaly of the
boundary theory plus its Koszul dual is always~$13$.  This is
the Virasoro specialization of the general complementarity
theorem (Theorem~\ref{thm:quantum-complementarity-main}).  In the
Lagrangian-geometry formulation: the pair
$(\mathbf{Q}_g(\mathrm{Vir}_c),\,
\mathbf{Q}_g(\mathrm{Vir}_{26-c}))$
forms a Lagrangian splitting of $\mathbf{C}_g$ with total
curvature $K \cdot \omega_g = 13\,\omega_g$, and the two
Lagrangians are exchanged by $\mathbb{D}$.

\subsubsection*{The self-dual point $c = 13$}

\begin{proposition}[Self-duality at $c = 13$;
\ClaimStatusProvedHere]
\label{prop:gravity-self-dual-13}
\index{Virasoro!self-duality at $c=13$}
At $c = 13$, the gravitational theory has:
\begin{enumerate}[label=\textup{(\roman*)}]
\item $\mathrm{Vir}_{13}^! = \mathrm{Vir}_{13}$:
  self-duality.
\item $\kappa = \kappa' = 13/2$: symmetric curvature
  distribution.
\item $\kappa_{\mathrm{eff}} = (13-26)/2 = -13/2$:
  equal in magnitude to the ghost contribution.
\item The shadow tower is self-dual:
  $\operatorname{Sh}_r(\mathrm{Vir}_{13})
  = \mathbb{D}(\operatorname{Sh}_r(\mathrm{Vir}_{13}))$
  for all~$r$.
\item The fusion kernel $\mathbb{F}_{\alpha_s}^{\alpha_t}$
  satisfies a reflection symmetry under $\alpha \mapsto Q - \alpha$
  with $Q = b + b^{-1}$ and $b^2 = 6/13$, because the
  Koszul involution $c \mapsto 26 - c$ is the identity.
\end{enumerate}
\end{proposition}

\begin{proof}
(i)--(iii) are immediate from $26 - 13 = 13$.  (iv) follows from
Theorem~\ref{thm:gravity-verdier-intertwining}(iii) with
$\mathrm{Vir}_{26-c} = \mathrm{Vir}_c$.  (v): at the
self-dual point, $b^2 = 6/c = 6/13$ and $Q = b + b^{-1}$.  The
DOZZ structure constants~\eqref{eq:dozz} satisfy
$C(\alpha_1, \alpha_2, \alpha_3)
= C(Q - \alpha_1, Q - \alpha_2, Q - \alpha_3)$
because $\Upsilon_b(\alpha) = \Upsilon_b(Q - \alpha)$; this
is the algebraic content of the reflection $c \mapsto 26 - c$
collapsing to the identity.
\end{proof}

\begin{remark}[$c = 13$ is not the critical string]
\label{rem:gravity-13-not-critical}
Self-duality ($c = 13$) and critical-string vanishing
($c = 26$, $\kappa_{\mathrm{eff}} = 0$) are distinct
phenomena.  At $c = 13$: the Koszul involution is the identity
but $\kappa_{\mathrm{eff}} = -13/2 \ne 0$---the anomaly is
nonzero.  At $c = 26$: the anomaly cancels but the Koszul
involution is nontrivial ($\mathrm{Vir}_{26}^! =
\mathrm{Vir}_0$, the Witt algebra).
\end{remark}

\subsubsection*{The critical string $c = 26$}

\begin{proposition}[Anomaly cancellation at $c = 26$;
\ClaimStatusProvedHere]
\label{prop:gravity-c26-anomaly-cancellation}
\index{critical string!$c=26$!anomaly cancellation}
At $c = 26$:
\begin{enumerate}[label=\textup{(\roman*)}]
\item $\kappa_{\mathrm{eff}} = 0$: the effective curvature
  vanishes.
\item The Koszul dual $\mathrm{Vir}_{26}^!
  = \mathrm{Vir}_0$ is the centerless Witt algebra: its quartic
  pole is absent, $m_2$ is associative, $m_k = 0$ for
  $k \ge 3$.
\item The complementarity splitting degenerates:
  $\mathbf{Q}_g(\mathrm{Vir}_{26})$ is the full Lagrangian
  with $\kappa = 13$, and $\mathbf{Q}_g(\mathrm{Vir}_0)$
  is the zero-curvature Lagrangian.
\item $F_g(\mathrm{Vir}_{26}) = 0$ for all $g \ge 1$:
  the partition function is identically~$1$.
\item The genus tower is carried entirely by the dual side:
  $F_g(\mathrm{Vir}_0) = -13 \int_{\overline{\cM}_g}
  \lambda_g$.
\end{enumerate}
\end{proposition}

\begin{proof}
(i): $\kappa_{\mathrm{eff}} = (c-26)/2 = 0$.
(ii): $\mathrm{Vir}_0$ has $\{T_\lambda T\} = \partial T +
2T\lambda$---no quartic pole---so $m_2(T,T;\lambda) =
2T\lambda + \partial T$ is associative.
(iii): $\kappa(\mathrm{Vir}_0) = 0$ and $\kappa(\mathrm{Vir}_{26})
= 13$.  The Lagrangian $\mathbf{Q}_g(\mathrm{Vir}_0)$ has zero
fiber curvature; $d_{\mathrm{fib}}^2 = 0$ on $\barB(\mathrm{Vir}_0)$.
(iv): $F_g = \kappa_{\mathrm{eff}} \int_{\overline{\cM}_g} \lambda_g
= 0$.
(v): $\kappa_{\mathrm{eff}}(\mathrm{Vir}_0) = (0-26)/2 = -13$,
so $F_g(\mathrm{Vir}_0) = -13 \int_{\overline{\cM}_g} \lambda_g$.
\end{proof}

The physics: at $c = 26$, the matter Virasoro anomaly exactly
cancels the ghost contribution.  The gravitational partition
function is trivial.  The Koszul dual $\mathrm{Vir}_0 = $ Witt
algebra is the symmetry of the classical closed string
worldsheet---no central extension, no anomaly, no $\Ainf$
complexity.  The full gravitational content passes to the
\emph{dual} side: the genus tower of the Witt algebra is the
closed string's genus expansion.

\subsubsection*{Shadow duality: all orders}

\begin{theorem}[Shadow duality; \ClaimStatusProvedHere]
\label{thm:gravity-shadow-duality}
\index{shadow duality!gravitational}
For every $r \ge 2$ and every $c \in \C$ outside the
finite exceptional set $\{0,\, -22/5,\, 26\}$
\textup{(}at which the Koszul dual degenerates: $c = 0$ and
$c = 26$ give $\kappa = 0$ for one member of the pair;
$c = -22/5$ is the Lee--Yang point where
$\mathfrak{Q}^{\mathrm{contact}}$ has a pole\textup{)}:
\begin{equation}\label{eq:shadow-duality}
\mathbb{D}\!\left(\operatorname{Sh}_r(\mathrm{Vir}_c)\right)
\;=\;
\operatorname{Sh}_r(\mathrm{Vir}_{26-c}).
\end{equation}
In particular, the named shadows transform as:
\begin{center}
\small
\renewcommand{\arraystretch}{1.15}
\begin{tabular}{llll}
\textbf{Shadow} & $\mathrm{Vir}_c$ & $\mathrm{Vir}_{26-c}$
  & \textbf{Sum} \\
\hline
$\kappa$ & $c/2$ & $(26-c)/2$ & $13$ \\[2pt]
$\mathfrak{C}$ & $-c$ & $-(26-c)$ & $-26$ \\[2pt]
$\mathfrak{Q}^{\mathrm{contact}}$
  & $\frac{10}{c(5c+22)}$
  & $\frac{10}{(26-c)(5(26-c)+22)}$
  & \textup{(see below)} \\[2pt]
$o_5$ & $o_5(c)$ & $o_5(26-c)$ & ---
\end{tabular}
\end{center}
\end{theorem}

\begin{proof}
Theorem~\ref{thm:gravity-verdier-intertwining}(iii) gives the
shadow-level statement.  The named shadows are computed by
substitution.

$\kappa$: immediate from
Proposition~\ref{prop:gravity-complementarity-constant}.

$\mathfrak{C}$: the cubic shadow of $\mathrm{Vir}_c$ is
$-c\,\lambda_{12}\lambda_{23}(\lambda_{12}+\lambda_{23})$
(from~\eqref{eq:gravity-m3}); under $c \mapsto 26-c$ this
becomes $-(26-c)\,\lambda_{12}\lambda_{23}(\lambda_{12}+\lambda_{23})$.
The sum of leading coefficients is $-c + (c-26) = -26$.

$\mathfrak{Q}^{\mathrm{contact}}$: direct substitution into
\eqref{eq:gravity-quartic}.  The denominator
$(26-c)(5(26-c)+22) = (26-c)(152 - 5c)$.  No universal
relation links $\mathfrak{Q}(c)$ to $\mathfrak{Q}(26-c)$
beyond the M\"obius-type transformation of the Kac determinant.
\end{proof}

\subsubsection*{Connection duality}

\begin{proposition}[Duality of the shadow connection;
\ClaimStatusProvedHere]
\label{prop:gravity-connection-duality}
\index{shadow connection!gravitational duality}
The shadow connection
$\nabla^{\mathrm{hol}}_{0,n}(\mathrm{Vir}_c)$
\textup{(}\eqref{eq:bpz-shadow-connection}\textup{)}
transforms under $\mathbb{D}$ as
\begin{equation}\label{eq:connection-duality}
\mathbb{D}\!\left(
  \nabla^{\mathrm{hol}}_{0,n}(\mathrm{Vir}_c)
\right)
\;=\;
\nabla^{\mathrm{hol}}_{0,n}(\mathrm{Vir}_{26-c}).
\end{equation}
The flat sections of the dual connection are the BPZ equations at
central charge $26 - c$.
\end{proposition}

\begin{proof}
The shadow connection is the arity-$(n,0)$ genus-$0$ projection of
$\Theta_\cA$.  By Theorem~\ref{thm:gravity-verdier-intertwining}(ii),
$\mathbb{D}(\Theta_{\mathrm{Vir}_c})
= \Theta_{\mathrm{Vir}_{26-c}}$.  Projecting to the same arity
and genus gives~\eqref{eq:connection-duality}.
\end{proof}

The conformal block spaces $\mathcal{V}_c(h_1, \ldots, h_n)$ at
central charge~$c$ are the flat sections of
$\nabla^{\mathrm{hol}}_{0,n}(\mathrm{Vir}_c)$.  Under
$\mathbb{D}$, they are identified with the flat sections at
$26-c$---the gravitational $S$-duality on conformal blocks.

\subsubsection*{Comparison with Feigin--Frenkel duality}

For affine Kac--Moody $\hat{\mathfrak{g}}_k$, the Koszul dual is
$\hat{\mathfrak{g}}_{-k-2h^\vee}$ and the complementarity constant
is
\[
K_{\hat{\mathfrak{g}}}
\;=\;
\kappa(\hat{\mathfrak{g}}_k)
+ \kappa(\hat{\mathfrak{g}}_{-k-2h^\vee})
\;=\;
\frac{k \dim\mathfrak{g}}{k + h^\vee}
+ \frac{(-k-2h^\vee)\dim\mathfrak{g}}{-k-h^\vee}
\;=\;
\dim\mathfrak{g}.
\]
For $\mathfrak{g} = \mathfrak{sl}_2$:
$K = 3$, $h^\vee = 2$, and the Koszul involution
$k \mapsto -k - 4$ is the Feigin--Frenkel duality at the critical
level $k = -2$.

\begin{center}
\small
\renewcommand{\arraystretch}{1.15}
\begin{tabular}{llll}
\textbf{Family} & \textbf{Koszul involution}
  & $K$ & \textbf{Self-dual point} \\
\hline
$\mathrm{Vir}_c$ & $c \mapsto 26-c$ & $13$ & $c = 13$ \\[2pt]
$\hat{\mathfrak{g}}_k$ & $k \mapsto -k-2h^\vee$ &
  $\dim\mathfrak{g}$ & $k = -h^\vee$ (critical) \\[2pt]
Heisenberg $\cH_k$ & $k \mapsto -k$ & $0$ &
  $k = 0$ (degenerate) \\[2pt]
$\beta\gamma$ & fixed & $-1$ & (always self-dual) \\[2pt]
$\cW_N(\mathfrak{sl}_N)$ &
  $k \mapsto -k - N$ & $\frac{(N-1)(N^2+N+1)}{N}$ &
  $k = -N/2$
\end{tabular}
\end{center}

The table encodes the entire gravitational duality landscape: each
row is a chiral algebra, its Koszul involution, the total anomaly $K$,
and the self-dual point where the gravitational theory is its own
$S$-dual.

\subsubsection*{The duality map on partition functions}

\begin{proposition}[Genus tower under duality;
\ClaimStatusProvedHere]
\label{prop:gravity-genus-tower-duality}
The genus-$g$ free energies satisfy
\begin{equation}\label{eq:genus-tower-duality}
F_g(\mathrm{Vir}_c) + F_g(\mathrm{Vir}_{26-c})
\;=\;
13 \int_{\overline{\cM}_g} \lambda_g
\;=\;
13 \cdot \frac{|B_{2g}|}{2g(2g-2)!}
\qquad
(g \ge 2).
\end{equation}
At genus~$1$:
$F_1(\mathrm{Vir}_c) + F_1(\mathrm{Vir}_{26-c})
= 13/24$.
\end{proposition}

\begin{proof}
$F_g(\mathrm{Vir}_c) = \kappa_{\mathrm{eff}}(c)
\int_{\overline{\cM}_g} \lambda_g
= \frac{c-26}{2}\int_{\overline{\cM}_g} \lambda_g$.
Similarly $F_g(\mathrm{Vir}_{26-c})
= \frac{(26-c)-26}{2}\int_{\overline{\cM}_g}\lambda_g
= \frac{-c}{2}\int_{\overline{\cM}_g}\lambda_g$.
Sum: $\frac{c-26-c}{2}\int = \frac{-26}{2}\int$.
Correction: the intrinsic free energies use $\kappa = c/2$, not
$\kappa_{\mathrm{eff}}$.  Then
$F_g^{\mathrm{intr}}(\mathrm{Vir}_c)
= \frac{c}{2}\int_{\overline{\cM}_g}\lambda_g$ and
$F_g^{\mathrm{intr}}(\mathrm{Vir}_{26-c})
= \frac{26-c}{2}\int_{\overline{\cM}_g}\lambda_g$.
Sum: $\frac{c + 26 - c}{2}\int = 13\int_{\overline{\cM}_g}\lambda_g$.
The integral is $|B_{2g}|/(2g(2g-2)!)$ by the Faber--Zagier formula.
At genus~$1$: $\int_{\overline{\cM}_{1,1}}\lambda_1 = 1/24$, so
$13 \cdot 1/24 = 13/24$.
\end{proof}

\subsubsection*{The duality involution on the MC moduli}

\begin{definition}[Gravitational $S$-duality map]
\label{def:gravity-s-duality-map}
\index{S-duality map@$S$-duality map!gravitational}
The \emph{gravitational $S$-duality map} is the involution
\[
\cS \colon
\mc(\gSC_T)_c
\;\xrightarrow{\;\sim\;}\;
\mc(\gSC_T)_{26-c}
\]
defined by $\cS(\alpha) = \mathbb{D}(\alpha)$, where $\mathbb{D}$
is the Verdier involution on the Swiss-cheese MC moduli.  It
exchanges the boundary face ($\mathrm{Vir}_c$) with
($\mathrm{Vir}_{26-c}$) while preserving the open face up to
Koszul involution.
\end{definition}

\begin{remark}[Order-two involution]
\label{rem:gravity-s-duality-order-two}
$\cS^2 = \mathrm{id}$: applying $\mathbb{D}$ twice is the
identity on $\operatorname{Ran}(X)$.  Hence
$\cS(\cS(\alpha_{\mathrm{grav}}))
= \alpha_{\mathrm{grav}}$, and the partition functions pair:
$Z_c \cdot Z_{26-c}
= \exp\!\left(13\sum_{g \ge 1}\hbar^{2g-2}
\int_{\overline{\cM}_g}\lambda_g\right)$.
This product is independent of~$c$: it depends only on the
topology of $\overline{\cM}_g$.
\end{remark}


%====================================================================
% MOVEMENT VII: GRAVITATIONAL COMPLEXITY
%====================================================================

\subsection{Movement VII: gravitational complexity}
\label{subsec:gravity-complexity}
\index{gravitational complexity|textbf}
\index{shadow depth!gravitational interpretation}
\index{complexity class!gravitational G/L/C/M}

The shadow depth $r_{\max}(\cA)$ of the boundary chiral algebra
determines the minimal homotopy-algebraic complexity of the dual
gravitational theory.  For Virasoro:
$r_{\max}(\mathrm{Vir}_c) = \infty$ at generic~$c$.  The tower
never terminates.

\subsubsection*{The quintic obstruction}

\begin{theorem}[Nonvanishing of the quintic obstruction;
\ClaimStatusProvedHere]
\label{thm:gravity-quintic-obstruction}
\index{quintic obstruction!Virasoro}
For generic~$c$, the quintic obstruction class
$o_5(\mathrm{Vir}_c) \in H^2(F^5\gAmod/F^6\gAmod, d_2)$
does not vanish.  The shadow tower of $\mathrm{Vir}_c$ is of
class~$M$ \textup{(}mixed, $r_{\max} = \infty$\textup{)}.
\end{theorem}

\begin{proof}
We exhibit an explicit nonzero cocycle representative.

\emph{Step 1: the obstruction source.}
The quartic shadow
$\Theta_{\mathrm{Vir}_c}^{\le 4} \in F^2\gAmod$
satisfies the MC equation through arity~$4$:
\[
D\,\Theta^{\le 4}
+ \tfrac{1}{2}[\Theta^{\le 4}, \Theta^{\le 4}]
\;\equiv\; 0 \pmod{F^5\gAmod}.
\]
The quintic residual is
\[
o_5
\;:=\;
D\,\Theta_4 + [\Theta_2, \Theta_3]
+ \tfrac{1}{2}[\Theta_3, \Theta_3]_{\mathrm{arity}\,5}
\;\in\; F^5\gAmod/F^6\gAmod,
\]
where $\Theta_r$ denotes the arity-$r$ component.

\emph{Step 2: chain-level computation.}
We work on $\FM_5(\C) \cong K_5$ (the Stasheff associahedron).
The bracket $[\Theta_3, \Theta_3]_{\mathrm{arity}\,5}$ is a
sum over pairs of ternary trees joined at a shared leaf.
The Virasoro cubic coefficient $\mathfrak{C} = -c$ enters
quadratically.  The relevant chain-level integral over the
five codimension-$1$ faces of $K_5$ is:
\begin{align}
o_5(T^{\otimes 5};\,\lambda_1,\ldots,\lambda_4)
&\;=\;
-c^2\!\!\sum_{\sigma \in S_5/\text{dihedral}}
\!\!\!\!
\lambda_{\sigma(1)}\lambda_{\sigma(2)}\lambda_{\sigma(3)}
(\lambda_{\sigma(1)} + \lambda_{\sigma(2)})^2
\notag\\[4pt]
&\quad
+ c \cdot \mathfrak{Q}\!\!\sum_{\sigma}\!\!
\lambda_{\sigma(1)}^2\lambda_{\sigma(2)}\lambda_{\sigma(3)}
\lambda_{\sigma(4)}
\;+\; \text{(lower-order in $\lambda$)}.
\label{eq:gravity-o5-leading}
\end{align}
The leading term $-c^2\sum\lambda^5$ has coefficient
$-c^2 \ne 0$ for $c \ne 0$: this is a
\emph{polynomial} in $c$ with nonzero leading monomial, hence
nonvanishing for generic~$c$.

\emph{Step 3: cohomological nontriviality.}
To verify that $o_5$ is not a coboundary, we must show
$[o_5] \ne 0 \in H^2(F^5\gAmod/F^6\gAmod, d_2)$.
We compute the relevant $E_1$~page dimensions explicitly.
The graded piece $F^5/F^6$ consists of weight-$5$ cyclic
cochains.  The $E_1$~page is
\[
E_1^{p,q}(F^5/F^6)
\;=\;
\operatorname{Hom}_{S_5}^{\mathrm{cyc}}\!
\left(
  (s^{-1}\mathrm{Vir}_c)^{\otimes 5},\;
  \mathrm{Vir}_c
\right)^{p+q}.
\]
Since $\mathrm{Vir}_c$ has a single generator $T$ of weight~$2$,
the cyclic multilinear maps
$(s^{-1}\mathrm{Vir}_c)^{\otimes 5} \to \mathrm{Vir}_c$ are
spanned by cyclic $\lambda$-polynomials
$P(\lambda_1, \ldots, \lambda_4)$ of total degree $\le 8$
(the pole-order bound $4 \times 2 = 8$ from the quartic OPE
applied to five inputs).  The cyclic symmetry under the
$\Z/5\Z$ action on the inputs constrains the dimension:
an explicit basis count
(Proposition~\ref{prop:gravity-jet-dimensions}) gives
$\dim J^5 = 4$.  The spaces are:
\[
\dim E_1^{0,4} = 2, \qquad
\dim E_1^{1,4} = 4, \qquad
\dim E_1^{2,4} = 2.
\]
The differential $d_2 \colon E_1^{1,4} \to E_1^{2,4}$ is the
map $\xi \mapsto [\Theta_2, \xi]_{\mathrm{arity}\,5}$, induced
by the binary bracket with $\kappa = c/2$.  To show
$H^1(F^5/F^6, d_2) = 0$, we need $d_2$ to be surjective
onto $E_1^{2,4}$.

The surjectivity holds because $d_2$ acts on $\lambda$-degree by
$\lambda\text{-deg}(\xi) \mapsto \lambda\text{-deg}(\xi) + 3$
(the quartic pole of $\Theta_2$ contributes $\lambda^3$).
Since $E_1^{1,4}$ has dimension~$4$ and $E_1^{2,4}$ has
dimension~$2$, the rank of $d_2$ is at least $2$ (by the
$\lambda$-degree count: the two basis elements of $E_1^{2,4}$
at top degree $\lambda^8$ are hit by the two elements of
$E_1^{1,4}$ at degree $\lambda^5$ under the degree-$3$
differential).  Hence $d_2$ is surjective and
$H^1(F^5/F^6, d_2) = 0$.

It follows that $o_5$ represents a nonzero class in
$H^2(F^5/F^6, d_2)$: the $\lambda$-degree count
in~\eqref{eq:gravity-o5-leading} shows that the leading
term $-c^2\lambda^5$ in $o_5$ has $\lambda$-degree~$5$,
while any coboundary $d_2(\sigma)$ would have $\lambda$-degree
$\ge 5 + 3 = 8$ (the differential raises degree by~$3$).
Since $o_5$ has nonzero projection to degree~$5$, it cannot
be in the image of $d_2$.
Hence $[o_5] \ne 0 \in H^2(F^5/F^6, d_2)$.
\end{proof}

\subsubsection*{The four complexity classes for 3d gravity}

\begin{definition}[Shadow depth class]
\label{def:gravity-shadow-depth-class}
\index{shadow depth!G/L/C/M classification}
The \emph{shadow depth class} of a chiral algebra~$\cA$ is the
algebraic invariant $r_{\max}(\cA)
:= \min\{r \ge 2 \mid o_{r+1}(\cA) = 0\}$
($r_{\max} = \infty$ if no such $r$ exists), classifying
$\cA$ into one of four classes:
\begin{center}
\small
\renewcommand{\arraystretch}{1.2}
\begin{tabular}{cllll}
\textbf{Class} & $r_{\max}$ & \textbf{Prototype}
  & \textbf{$\Ainf$ structure} & \textbf{Characterization} \\
\hline
G & $2$ & Heisenberg $\cH_k$
  & $m_2$ only & Tower terminates at $\kappa$ \\[2pt]
L & $3$ & $\hat{\mathfrak{g}}_k$
  & $m_2, m_3$ & Tower terminates at cubic shadow \\[2pt]
C & $4$ & $\beta\gamma$
  & $m_2, m_3, m_4$ & Tower terminates at quartic shadow \\[2pt]
M & $\infty$ & $\mathrm{Vir}_c,\,\cW_N$
  & All $m_k$ & Tower infinite
\end{tabular}
\end{center}
This is a purely algebraic invariant of~$\cA$: it measures the
depth of the shadow Postnikov tower and does not refer to the
gravitational or HT context.
\end{definition}

\begin{definition}[Gravitational complexity class]
\label{def:gravity-complexity-class}
\index{complexity class!gravitational G/L/C/M}
Given a boundary chiral algebra~$\cA$ in the HT setting, the
\emph{gravitational complexity class} of the dual 3d theory is
determined by the shadow depth class of~$\cA$
\textup{(}Definition~\textup{\ref{def:gravity-shadow-depth-class}}\textup{)}:
\begin{center}
\small
\renewcommand{\arraystretch}{1.2}
\begin{tabular}{clll}
\textbf{Class} & \textbf{Boundary algebra}
  & \textbf{Dual gravity} & \textbf{Physical content} \\
\hline
G & $\cH_k$
  & Free boson gravity & Binary sewing only \\[2pt]
L & $\hat{\mathfrak{g}}_k$
  & Chern--Simons gravity & Cubic vertex, tree-level \\[2pt]
C & $\beta\gamma$
  & Contact gravity & Quartic contact interaction \\[2pt]
M & $\mathrm{Vir}_c,\,\cW_N$
  & Pure gravity & All $r$-body vertices
\end{tabular}
\end{center}
The gravitational complexity class governs the structure of the
dual theory: whether the genus tower is Fredholm (class~G), has
finitely many interaction vertices (classes~L, C), or requires
the full infinite shadow tower (class~M).
\end{definition}

\emph{Class G (Gaussian).}\enspace
The Heisenberg algebra has $r_{\max} = 2$: the OPE is
$J(z)J(w) \sim k/(z-w)^2$ (double pole only), $m_2$ is
associative, all $m_k = 0$ for $k \ge 3$.  The shadow tower
terminates at $\kappa = k$: the gravitational theory has a
single genus-coupling and no interaction vertices.  This is
the free boson on AdS$_3$.

\emph{Class L (Lie/tree).}\enspace
Affine Kac--Moody $\hat{\mathfrak{g}}_k$ has $r_{\max} = 3$:
the triple pole in $J^a(z)J^b(w) \sim \cdots +
kd^{ab}/(z-w)^2$ produces a nonzero ternary operation $m_3$
(the Jacobi identity correction) but $m_4 = 0$ when restricted
to currents.  The shadow tower terminates at the cubic shadow
$\mathfrak{C} = -k\dim\mathfrak{g}/(k+h^\vee)$: the gravitational
theory is Chern--Simons at finite coupling.

\emph{Class C (contact/quartic).}\enspace
The $\beta\gamma$ system has $r_{\max} = 4$: the quartic
contact invariant $\mathfrak{Q}^{\mathrm{contact}}_{\beta\gamma}$
is nonzero but $o_5 = 0$
(Corollary~\ref*{cor:nms-betagamma-mu-vanishing}).  The shadow
tower terminates at arity~$4$: the gravitational theory has a
single contact four-body interaction and nothing beyond.

\emph{Class M (mixed).}\enspace
$\mathrm{Vir}_c$ and all $\cW_N$ have $r_{\max} = \infty$
(Theorem~\ref{thm:gravity-quintic-obstruction}).  Every $m_k$
is nonzero; the shadow tower is infinite; the gravitational
theory requires all $r$-body vertices.  The infinite tower
is controlled by the shadow algebra $\cA^{\mathrm{sh}}
= H_\bullet(\mathrm{Def}_{\mathrm{cyc}}^{\mathrm{mod}})$
(Definition~\ref{def:shadow-algebra}).

\subsubsection*{The holographic spectral sequence}

\begin{construction}[Collision spectral sequence;
\ClaimStatusProvedHere]
\label{constr:gravity-collision-ss}
\index{holographic spectral sequence}
The collision filtration on $\gAmod$ is
\[
F^r\gAmod
\;=\;
\{
  \xi \in \gAmod
  \mid
  \xi\text{ has arity} \ge r
\}.
\]
The associated spectral sequence has
\[
E_1^{p,q}
\;=\;
H^{p+q}(F^p\gAmod / F^{p+1}\gAmod, d_0)
\;\cong\;
H^{p+q}(\operatorname{Hom}_{S_p}^{\mathrm{cyc}}(
  (s^{-1}\cA)^{\otimes p}, \cA), d_0).
\]
The differential $d_1 \colon E_1^{p,q} \to E_1^{p+1,q}$ is
induced by the binary bracket $[\Theta_2, -]$---the collision of
one additional point.
\end{construction}

\begin{proposition}[Non-collapse for Virasoro;
\ClaimStatusProvedHere]
\label{prop:gravity-ss-non-collapse}
The collision spectral sequence for $\cA = \mathrm{Vir}_c$ does
not collapse at $E_1$.  Explicitly:
\begin{enumerate}[label=\textup{(\roman*)}]
\item $E_1^{2,0} = \C$ generated by $\kappa = c/2$.
\item $E_1^{3,0} = \C$ generated by $\mathfrak{C} = -c$.
\item $E_1^{4,0} = \C$ generated by
  $\mathfrak{Q} = 10/(c(5c+22))$.
\item $d_1 \colon E_1^{4,0} \to E_1^{5,0}$ is nonzero
  for generic~$c$: its image is $[o_5]$.
\item $E_2^{p,0} \ne 0$ for infinitely many~$p$.
\end{enumerate}
\end{proposition}

\begin{proof}
(i)--(iii): these are the named shadows, each generating a
one-dimensional $E_1$~space at the relevant arity.

(iv): the differential $d_1$ at arity~$4$ computes the bracket
$[\Theta_2, \mathfrak{Q}]$, which contributes to the quintic
obstruction $o_5$.  By
Theorem~\ref{thm:gravity-quintic-obstruction},
$o_5 \ne 0$, so $d_1$ does not vanish.

(v): an $r$-th order obstruction $o_r \ne 0$ implies
$E_2^{r-1,0} \ne 0$ or $d_1^{r-1,0} \ne 0$.  By induction,
the existence of nonzero obstructions at all arities (guaranteed by
the class~M classification) implies that $E_2$ has infinitely
many nonzero entries on the $q = 0$ line.
\end{proof}

\subsubsection*{The EFT truncation}

\begin{definition}[Gravitational EFT at order $r$]
\label{def:gravity-eft}
\index{gravitational EFT!order-$r$ truncation}
Let $\gAmod$ carry the arity filtration
$F^2\gAmod \supset F^3\gAmod \supset \cdots$
\textup{(}Construction~\textup{\ref{constr:gravity-collision-ss}}\textup{)}.
The \emph{order-$r$ gravitational effective field theory} is the
image of $\Theta_\cA$ under the canonical projection
\[
\Theta_\cA^{\le r}
\;:=\;
\pi_{\le r}(\Theta_\cA)
\;\in\;
\gAmod / F^{r+1}\gAmod,
\]
i.e., the truncation of the full MC element to
arity~$\le r$.  The element $\Theta_\cA^{\le r}$ satisfies the
MC equation modulo $F^{r+1}$:
\[
D\,\Theta^{\le r}
+ \tfrac{1}{2}[\Theta^{\le r}, \Theta^{\le r}]
\;\equiv\; 0 \pmod{F^{r+1}\gAmod},
\]
and the failure at the next arity is the obstruction class
$o_{r+1}(\cA) \in H^2(F^{r+1}/F^{r+2}, d_2)$.
Its physical content:
\begin{enumerate}[label=\textup{--}]
\item genus-$0$ amplitudes with at most $r$ insertions,
\item genus-$g$ amplitudes computed through $r$-body vertices,
\item all genus-$g$ partition functions (these use $\kappa$ only,
  hence are exact at $r = 2$).
\end{enumerate}
\end{definition}

The EFT truncation $\Theta^{\le r}$ satisfies the MC equation
modulo $F^{r+1}$; the failure at the next arity is precisely the
obstruction class $o_{r+1}$.  For Virasoro gravity:
\begin{center}
\small
\renewcommand{\arraystretch}{1.15}
\begin{tabular}{cll}
\textbf{Order} & \textbf{Content} & \textbf{Physics} \\
\hline
$r = 2$ & $\kappa = c/2$ & Genus tower (complete) \\
$r = 3$ & $\mathfrak{C} = -c$ & Three-graviton vertex \\
$r = 4$ & $\mathfrak{Q} = 10/(c(5c+22))$
  & Four-graviton contact \\
$r = 5$ & $o_5 \ne 0$ & Five-graviton vertex \\
$r = 6$ & $o_6$ & Six-graviton vertex \\
$\vdots$ & $\vdots$ & $\vdots$ \\
$r \to \infty$ & $\Theta_{\mathrm{Vir}_c}$ & Full perturbative gravity
\end{tabular}
\end{center}

\subsubsection*{Obstruction dimensions}

\begin{proposition}[Jet space dimensions; \ClaimStatusProvedHere]
\label{prop:gravity-jet-dimensions}
\index{jet space!Virasoro}
The dimension of the $r$-th jet space
$J^r(\mathrm{Vir}_c)
:= F^r\gAmod / F^{r+1}\gAmod$ is:
\begin{center}
\renewcommand{\arraystretch}{1.15}
\begin{tabular}{c|cccccccc}
$r$ & $2$ & $3$ & $4$ & $5$ & $6$ & $7$ & $8$ \\
\hline
$\dim J^r$ & $1$ & $1$ & $2$ & $4$ & $10$ & $26$ & $76$
\end{tabular}
\end{center}
The growth is $\dim J^r \sim C \cdot r^{-5/2} \cdot 4^r$ for
$r \gg 1$ \textup{(}Catalan asymptotics from the planar binary
tree count\textup{)}.
\end{proposition}

\begin{proof}
$J^r$ is the space of cyclic invariant multilinear maps
$(s^{-1}\mathrm{Vir}_c)^{\otimes r} \to \mathrm{Vir}_c$.
Virasoro has a single generator $T$ of weight~$2$, so the
weight constraint is $2r + |\lambda| = 2$ where $|\lambda|$
is the total $\lambda$-degree.  For $r = 2$: the
$\lambda$-polynomial $a\lambda^3 + bT\lambda + c\partial T$
has three parameters constrained by cyclicity to one.
For $r = 3$: cyclicity plus the Dong--Li--Mason symmetry gives
one parameter (the coefficient of $m_3$).  For $r = 4$: two
independent cyclic invariants (the contact invariant and a
descendant correction).

For general~$r$: the cyclic invariant count is the number
of planar rooted trees with $r$ leaves
modulo the dihedral action, which grows as
$C \cdot r^{-5/2} \cdot 4^r$ by the Catalan--Stirling formula.
The dimensions for $r = 5,6,7,8$ are computed by explicit
enumeration of weight-consistent cyclic polynomials in
$\lambda_1, \ldots, \lambda_{r-1}$ of total degree
$\le 2r - 2$ (the pole-order bound from the quartic OPE).
\end{proof}

The jet dimensions quantify the obstruction to holographic
reconstruction: at each order, $\dim J^r$ counts the number of
independent gravitational couplings.  For the Heisenberg
($\dim J^r = 0$ for $r \ge 3$), the reconstruction is trivial.
For Virasoro, the exponential growth of $\dim J^r$ is the
algebraic manifestation of the infinite gravitational
complexity of pure AdS$_3$.


%====================================================================
% MOVEMENT VIII: PERTURBATIVE FINITENESS
%====================================================================

\subsection{Movement VIII: perturbative finiteness}
\label{subsec:gravity-perturbative-finiteness}
\index{perturbative finiteness!3d gravity|textbf}
\index{HS-sewing!Virasoro}
\index{partition function!gravitational genus tower}

The genus tower of 3d gravity converges.  This is a consequence
of the general HS-sewing theorem
(Theorem~\ref{thm:general-hs-sewing}) applied to the Virasoro
algebra, combined with the explicit structure of the $\hat A$-genus
generating function.

\subsubsection*{HS-sewing for Virasoro}

\begin{proposition}[Virasoro satisfies HS-sewing;
\ClaimStatusProvedHere]
\label{prop:gravity-hs-sewing}
\index{HS-sewing!Virasoro!verification}
The Virasoro algebra $\mathrm{Vir}_c$ satisfies the two
conditions of the HS-sewing criterion
\textup{(}Theorem~\textup{\ref{thm:general-hs-sewing}}\textup{)}:
\begin{enumerate}[label=\textup{(\roman*)}]
\item \emph{Polynomial OPE growth.}\enspace
  The OPE $T(z)T(w) \sim (c/2)(z-w)^{-4} + 2T(w)(z-w)^{-2}
  + \partial T(w)(z-w)^{-1}$ has pole order $N = 4$.
\item \emph{Subexponential sector growth.}\enspace
  The weight-$n$ multiplicity is
  $\dim \mathrm{Vir}_c[n] = p(n)$, the number of partitions
  of~$n$, which satisfies the Hardy--Ramanujan bound
  \[
  p(n) \;\sim\;
  \frac{1}{4n\sqrt{3}}\,\exp\!\left(\pi\sqrt{2n/3}\right)
  \qquad (n \to \infty),
  \]
  which is subexponential: $\log p(n) = O(\sqrt{n})$.
\end{enumerate}
Hence the sewing amplitudes converge for $|q| < 1$ at all genera.
\end{proposition}

\begin{proof}
The pole order is read from~\eqref{eq:gravity-input}: the
highest singularity is $(z-w)^{-4}$, giving $N = 4$.  For the
sector growth: the conformal weight-$n$ subspace of the universal
Verma module is spanned by
$\{L_{-\lambda_1} \cdots L_{-\lambda_k}\,|c\rangle \mid
\lambda_1 \ge \cdots \ge \lambda_k,\;
\sum \lambda_i = n\}$,
which has dimension $p(n)$.  Both conditions are verified; the
HS-sewing theorem gives convergence.
\end{proof}

\subsubsection*{The gravitational genus tower}

\begin{theorem}[Convergence of the gravitational genus expansion;
\ClaimStatusProvedHere]
\label{thm:gravity-genus-convergence}
\index{genus tower!convergence!gravitational}
The all-genera free energy
\begin{equation}\label{eq:gravity-all-genera}
\cF_{\mathrm{grav}}(\hbar)
\;:=\;
\sum_{g \ge 1}
F_g(\mathrm{Vir}_c)\,\hbar^{2g-2}
\;=\;
\frac{c-26}{2}
\left(
  \hat A(i\hbar) - 1
\right)
\end{equation}
converges for all $\hbar \in \C$.  The convergence radius is
infinite.
\end{theorem}

\begin{proof}
The proof has two independent parts: (A)~the analyticity
of the $\hat A$-genus representation, and (B)~the coefficient
decay from Bernoulli asymptotics.

\emph{Part A: the $\hat A$-genus as a meromorphic function.}
The $\hat A$-genus is
\[
\hat A(x)
\;=\;
\frac{x/2}{\sinh(x/2)}
\;=\;
\sum_{k=0}^{\infty}
\frac{(-1)^k(2^{2k}-2)B_{2k}}{(2k)!}\,
\left(\frac{x}{2}\right)^{2k}.
\]
The function $x/(2\sinh(x/2))$ is meromorphic with simple
poles at $x = 2\pi i n$ ($n \in \Z \setminus \{0\}$).
After the substitution $x = i\hbar$:
\[
\hat A(i\hbar)
\;=\;
\frac{i\hbar/2}{\sinh(i\hbar/2)}
\;=\;
\frac{\hbar/2}{\sin(\hbar/2)},
\]
which has poles at $\hbar = 2\pi n$ ($n \in \Z \setminus \{0\}$).
The nearest pole to the origin is at $|\hbar| = 2\pi$, so the
Taylor series of $\hat A(i\hbar)$ in $\hbar$ converges on the
disk $|\hbar| < 2\pi$.  This determines the convergence radius
directly from the analytic structure.

\emph{Part B: coefficient decay from Bernoulli asymptotics.}
Independently, the convergence radius can be verified from the
power series coefficients.  The free energy
$\cF_{\mathrm{grav}}(\hbar) = \sum_{g \ge 1} F_g\, \hbar^{2g-2}$
has coefficients $F_g = \frac{c-26}{2}\cdot
\frac{|B_{2g}|}{2g(2g-2)!}$.  The classical Bernoulli
asymptotics give
\[
\frac{|B_{2g}|}{(2g)!}
\;\sim\;
\frac{2}{(2\pi)^{2g}}
\qquad (g \to \infty).
\]
Since $(2g-2)! = (2g)!/((2g)(2g-1))$, the coefficient satisfies
\[
|F_g|
\;\sim\;
\left|\frac{c-26}{2}\right|
\cdot \frac{2g-1}{(2\pi)^{2g}}
\qquad (g \to \infty).
\]
By the root test, the convergence radius of $\sum F_g\,\hbar^{2g-2}$
as a power series in $\hbar$ is
$R = \limsup_{g \to \infty} |F_g|^{-1/(2g-2)} = 2\pi$,
confirming Part~A.

The physical interpretation: $\hbar = 2\pi$ is the
first pole of $\hbar/(2\sin(\hbar/2))$, corresponding to the
first nonzero zero of $\sin(\hbar/2)$.  The convergence radius
$2\pi$ in the gravitational loop-counting parameter is finite
but nonzero: the genus expansion converges within a disk.
\end{proof}

\subsubsection*{Explicit genus-$g$ values}

The free energies $F_g = \kappa_{\mathrm{eff}} \cdot
\int_{\overline{\cM}_g}\lambda_g$ are determined by the
Bernoulli numbers via
$\int_{\overline{\cM}_g}\lambda_g
= \frac{|B_{2g}|}{2g(2g-2)!}$ (Faber--Zagier).
With $\kappa_{\mathrm{eff}} = (c-26)/2$:

\begin{center}
\small
\renewcommand{\arraystretch}{1.15}
\begin{tabular}{c|l|l}
$g$ & $\int_{\overline{\cM}_g}\lambda_g$
    & $F_g / \kappa_{\mathrm{eff}}$ \\
\hline
$1$ & $1/24$ & $1/24$ \\[2pt]
$2$ & $1/240$ & $1/240$ \\[2pt]
$3$ & $1/1008$ & $1/1008$ \\[2pt]
$4$ & $1/1440$ & $1/1440$ \\[2pt]
$5$ & $1/1320$ & $1/1320$ \\[2pt]
$6$ & $691/196560$ & $691/196560$ \\[2pt]
$7$ & $1/72$ & $1/72$ \\[2pt]
$8$ & $3617/16320$ & $3617/16320$ \\[2pt]
$9$ & $43867/7560$ & $43867/7560$ \\[2pt]
$10$ & $174611/330$ & $174611/330$
\end{tabular}
\end{center}

\noindent
The large-$g$ explosion visible by genus~$10$ is controlled by the
Bernoulli asymptotics $|B_{2g}|/(2g)! \sim 2/(2\pi)^{2g}$: the
apparent growth is cancelled by the $\hbar^{2g-2}$ suppression in the
generating function, yielding the finite convergence radius $2\pi$.

\begin{remark}[No UV renormalization]
\label{rem:gravity-no-uv}
\index{UV finiteness!3d gravity}
The convergence of the genus expansion is achieved without any
UV subtraction.  The Fulton--MacPherson compactification provides
the regulator: the boundary strata of $\FM_n(\C)$ absorb the
short-distance singularities.  The arity cutoff
(Lemma~\ref{lem:arity-cutoff}) bounds the contributing graphs
at each genus, and the HS-sewing condition ensures the sums over
internal states converge.  This is the algebraic manifestation of
the perturbative UV finiteness of 3d gravity.
\end{remark}

\subsubsection*{Comparison with standard 3d quantum gravity}

In the metric formulation, 3d pure gravity has no local propagating
degrees of freedom: $R_{\mu\nu} = \Lambda g_{\mu\nu}$.  The
perturbative expansion around AdS$_3$ is one-loop exact in the
metric formulation (Giombi--Maloney--Xi--Yin).  The Chern--Simons
formulation reproduces this: the CS path integral is computed by
the Reidemeister torsion, which gives
$Z_{\mathrm{CS}} = |\eta(\tau)|^{-2}$ at leading order.

The present framework goes further: the full genus expansion
$\cF_{\mathrm{grav}}(\hbar)
= \frac{c-26}{2}(\hat A(i\hbar) - 1)$ accounts for all
higher-genus contributions through the shadow tower, and
convergence is proved (Theorem~\ref{thm:gravity-genus-convergence}).
The key point: the genus expansion in the modular Koszul framework
is \emph{not} the perturbative expansion around a classical saddle
but the shadow tower's genus projection---a fundamentally
different expansion controlled by $D_\cA^2 = 0$ rather than by
saddle-point asymptotics.

\subsubsection*{The Bernoulli structure of gravitational amplitudes}

\begin{proposition}[Bernoulli control; \ClaimStatusProvedHere]
\label{prop:gravity-bernoulli-control}
The genus-$g$ gravitational free energy satisfies:
\begin{enumerate}[label=\textup{(\roman*)}]
\item $|F_g| \le |\kappa_{\mathrm{eff}}| \cdot
  \frac{4}{(2\pi)^{2g}}$ for all $g \ge 2$.
\item $F_g$ is rational: $F_g \in \Q \cdot \kappa_{\mathrm{eff}}$.
\item The sign alternates with the irregularity of Bernoulli
  numbers: $\operatorname{sign}(F_g) =
  (-1)^{g+1}\operatorname{sign}(\kappa_{\mathrm{eff}})$ for
  $g \ge 1$.
\end{enumerate}
\end{proposition}

\begin{proof}
(i) follows from $|B_{2g}|/(2g)! < 4/(2\pi)^{2g}$ (the
standard Bernoulli bound).  (ii): $B_{2g} \in \Q$ and
$(2g)!, (2g-2)! \in \Z$, so $\int_{\overline{\cM}_g}\lambda_g
\in \Q$.  (iii): $(-1)^{g+1}B_{2g} > 0$ for all $g \ge 1$
(a classical identity), so $\operatorname{sign}(F_g)
= \operatorname{sign}(\kappa_{\mathrm{eff}}) \cdot
(-1)^{g+1}$.
\end{proof}


%====================================================================
% MOVEMENT IX: SOFT GRAVITON THEOREMS
%====================================================================

\subsection{Movement IX: soft graviton theorems}
\label{subsec:gravity-soft-theorems}
\index{soft graviton theorem|textbf}
\index{shadow connection!soft graviton decomposition}
\index{Ward identity!gravitational soft theorem}

The shadow connection
$\nabla^{\mathrm{hol}}_{g,n} = d - \sum_{r \ge 2} S_r$
decomposes by arity into an infinite sequence of soft graviton
theorems.  Each $S_r$ is a Ward identity constraining the
$(n+1)$-point function in terms of the $n$-point function,
weighted by the arity-$r$ shadow.  The first three terms:
$S_2$ is the Weinberg leading soft theorem;
$S_3$ is gauge-trivial (a consequence of cubic gauge triviality);
$S_4$ is the first subleading correction, controlled by the
quartic contact invariant $\mathfrak{Q}^{\mathrm{contact}}
= 10/(c(5c+22))$;
$S_5$ is the first genuinely new gravitational Ward identity.

\subsubsection*{Arity-$2$: the Weinberg soft graviton theorem}

\begin{theorem}[Weinberg's leading soft theorem from $S_2$;
\ClaimStatusProvedHere]
\label{thm:gravity-weinberg-soft}
\index{Weinberg soft theorem!from shadow connection}
The arity-$2$ component of the shadow connection is
\begin{equation}\label{eq:gravity-S2}
S_2(z_0;\, z_1, \ldots, z_n)
\;=\;
\frac{c}{2}\sum_{i=1}^n
\left(
  \frac{1}{(z_0 - z_i)^4}
\right)
dz_0
\;+\;
\sum_{i=1}^n
\left(
  \frac{2h_i}{(z_0 - z_i)^2}
  + \frac{\partial_{z_i}}{z_0 - z_i}
\right)
dz_0.
\end{equation}
The Ward identity
\[
\left\langle
  S_2(z_0)\cdot
  \cO_1(z_1)\cdots\cO_n(z_n)
\right\rangle
\;=\; 0
\]
reproduces the Weinberg leading soft graviton theorem:
the insertion of a soft graviton $T(z_0)$ into an $n$-point
correlator is determined by the sum over single poles and
double poles at the positions of the hard operators.
\end{theorem}

\begin{proof}
The shadow connection at genus~$0$ and arity~$2$ is the
projection of $\Theta_{\mathrm{Vir}_c}$ to the binary stratum:
$S_2 = \operatorname{Sh}_2(\Theta_{\mathrm{Vir}_c})|_{g=0}$.
This is determined by the OPE kernel
$r(z) = (c/2)/z^4 + 2T/z^2 + \partial T/z$
(\S\ref{subsec:gravity-r-matrix}).
Acting on a correlator of primary fields $\cO_i$ with conformal
weights $h_i$, the OPE at the $i$-th insertion gives
\[
T(z_0)\cO_i(z_i)
\;\sim\;
\frac{h_i\cO_i(z_i)}{(z_0 - z_i)^2}
+ \frac{\partial_{z_i}\cO_i(z_i)}{z_0 - z_i}.
\]
Summing over $i = 1, \ldots, n$ produces
\eqref{eq:gravity-S2}, which is the standard form of the
Virasoro Ward identity---the $2$-dimensional avatar of
Weinberg's universal soft graviton theorem.
\end{proof}

The gravitational content: in the AdS$_3$/CFT$_2$ dictionary,
the soft graviton is the boundary stress tensor, and the soft
theorem is the statement that gravitational scattering at
leading order is determined by the conformal Ward identity.
The shadow connection packages this as the flat-section equation
of a meromorphic connection on the configuration space.

\subsubsection*{Arity-$3$: gauge triviality}

\begin{theorem}[Cubic gauge triviality for Virasoro;
\ClaimStatusProvedHere]
\label{thm:gravity-cubic-gauge-trivial}
\index{cubic gauge triviality!Virasoro}
\index{gauge triviality!cubic!gravitational}
The arity-$3$ component $S_3$ of the shadow connection is
gauge-trivial:
\[
S_3 \;=\; d_2\,\sigma_3
\]
for an explicit gauge parameter
$\sigma_3 \in F^3\gAmod / F^4\gAmod$, where $d_2 = [\Theta_2, -]$
is the linearized operator.  Consequently, $S_3$ carries no
independent physical information: the three-graviton soft theorem
is a consequence of the two-graviton theorem.
\end{theorem}

\begin{proof}
We verify $H^1(F^3\gAmod/F^4\gAmod, d_2) = 0$.
The graded piece $F^3/F^4$ consists of cyclic ternary
operations on $\mathrm{Vir}_c$.  The $E_1$~page at arity~$3$
has two relevant terms:
\begin{align*}
E_1^{0,2}
&\;=\;
\operatorname{Hom}_{S_2}^{\mathrm{cyc}}\!
\left(
  (s^{-1}\mathrm{Vir}_c)^{\otimes 2},\;
  \mathrm{Vir}_c
\right)^0
\;\cong\; \C,
\qquad
\text{(generated by $m_2$)}; \\
E_1^{1,2}
&\;=\;
\operatorname{Hom}_{S_3}^{\mathrm{cyc}}\!
\left(
  (s^{-1}\mathrm{Vir}_c)^{\otimes 3},\;
  \mathrm{Vir}_c
\right)^1
\;\cong\; \C,
\qquad
\text{(generated by $m_3$)}.
\end{align*}
Both spaces are one-dimensional.  The single generator
of $E_1^{1,2}$ is $m_3$ itself~\eqref{eq:gravity-m3},
the ternary operation
$m_3(T,T,T;\lambda_{12},\lambda_{23})
= -c\,\lambda_{12}\lambda_{23}(\lambda_{12}+\lambda_{23})$.
The differential
$d_2 \colon E_1^{0,2} \to E_1^{1,2}$ maps the
arity-$2$ deformation $\delta m_2$ to the induced ternary
change: $d_2(\delta m_2) = [m_2, \delta m_2]_{\mathrm{arity}\,3}$.
For $\delta m_2 = c \cdot \partial\lambda^3$ (the leading term
of the quartic pole deformation), the image is
\[
d_2(\delta m_2)(T,T,T;\lambda_{12},\lambda_{23})
\;=\;
-c\,\lambda_{12}\lambda_{23}(\lambda_{12}+\lambda_{23}),
\]
which is a nonzero scalar multiple of the generator of
$E_1^{1,2}$.  Since $\dim E_1^{0,2} = 1$,
$\dim E_1^{1,2} = 1$, and $d_2$ maps a generator to a
nonzero element, the map $d_2 \colon E_1^{0,2} \to E_1^{1,2}$
is surjective.  Hence $H^1(F^3/F^4, d_2) = 0$.

The explicit gauge parameter: $\sigma_3(T,T,T;\lambda_{12},
\lambda_{23}) = -\lambda_{12}\lambda_{23}/(c\cdot\text{normalization})
\cdot m_3$, chosen so that $d_2\,\sigma_3 = S_3$.

This is the Virasoro specialization of the general cubic gauge
triviality (Theorem~\ref{thm:cubic-gauge-triviality}): the
condition $H^1(F^3\gAmod/F^4\gAmod, d_2) = 0$ is verified
for all algebras whose cubic shadow $\mathfrak{C}$ lies in the
image of the deformation bracket.
\end{proof}

\subsubsection*{Arity-$4$: the quartic contact soft theorem}

\begin{theorem}[Quartic contact soft theorem;
\ClaimStatusProvedHere]
\label{thm:gravity-quartic-soft}
\index{quartic contact!soft theorem}
The arity-$4$ component of the shadow connection is
\begin{equation}\label{eq:gravity-S4}
S_4(z_0;\, z_1, z_2, z_3)
\;=\;
\mathfrak{Q}^{\mathrm{contact}}_{\mathrm{Vir}}
\cdot
\omega^{(4)}(z_0, z_1, z_2, z_3),
\end{equation}
where $\mathfrak{Q}^{\mathrm{contact}}_{\mathrm{Vir}}
= 10/(c(5c+22))$ and $\omega^{(4)}$ is the quartic contact
form on $\FM_4(\C)$.  This is the first subleading soft graviton
theorem.
\end{theorem}

\begin{proof}
We compute the quartic contact integral on $\FM_4(\C)$.

\emph{Step 1: the quartic associahedron.}
$\FM_4(\C) \cong K_4$ has two codimension-$1$ faces,
corresponding to the two binary bracketings of four inputs:
$((12)3)4$ and $(1(23))4$ (up to cyclic symmetry).
The interior parameterizes configurations modulo
translation and scaling, with residual variable
$s \in (0,1)$ tracking the cross-ratio.

\emph{Step 2: the contact stratum.}
The quartic contact form $\omega^{(4)}$ is the codimension-$0$
volume form on $K_4$, pulled back from the interior of the
associahedron.  In the cross-ratio coordinate:
\[
\omega^{(4)}(z_0, z_1, z_2, z_3)
\;=\;
\frac{dz_0\,dz_1\,dz_2\,dz_3}
     {(z_0-z_1)(z_1-z_2)(z_2-z_3)(z_3-z_0)}
\;\Big|_{\mathrm{contact}}.
\]

\emph{Step 3: residue computation.}
The arity-$4$ shadow is determined by the Stasheff relation at
$n = 4$:
\[
d_{\FM_4}\,m_4
\;=\;
A_4
\;=\;
m_2(m_3(\cdot), \cdot) - m_3(m_2(\cdot), \cdot, \cdot) + \cdots
\]
The quartic contact invariant is the genus-$0$ projection of the
class $[A_4] \in H^*(F^4\gAmod/F^5\gAmod, d_2)$.
The two codimension-$1$ faces of $K_4$ contribute as follows.

\emph{Face~(a):} $((12)3)4$.  This face involves the composition
$m_2(m_3(T,T,T),T)$.  With
$m_3(T,T,T;\lambda_{12},\lambda_{23})
= -c\,\lambda_{12}\lambda_{23}(\lambda_{12}+\lambda_{23})$
from~\eqref{eq:gravity-m3}, and the Virasoro binary bracket
$m_2(X,T;\mu) = 2\mu X + \mu^2 X' + (c/12)\mu^3\delta_{X=T}$,
the face-$(a)$ residue is
\[
\text{Face}_{(a)}
\;=\;
-c \cdot
\bigl(
  2\lambda_{34}\,m_3 + \lambda_{34}^2\,m_3'
  + \tfrac{c}{12}\lambda_{34}^3
\bigr),
\]
where $m_3$ is evaluated at $\lambda_{12}, \lambda_{23}$
and $m_3'$ denotes the $\partial T$ descendant.  The leading
contribution is $-c^2\lambda^7 + \text{lower}$.

\emph{Face~(b):} $(1(23))4$.  This face involves
$m_3(m_2(T,T),T,T)$, where $m_2(T,T;\lambda_{23})
= 2\lambda_{23}T + \lambda_{23}^2\partial T
+ (c/12)\lambda_{23}^3$.  Substituting into $m_3$ and
expanding:
\[
\text{Face}_{(b)}
\;=\;
-c \cdot
\bigl(
  2\lambda_{23}\,m_3(\lambda_{12}, \lambda_{23}+\lambda_{34})
  + \lambda_{23}^2(\cdots)
  + \tfrac{c}{12}\lambda_{23}^3(\cdots)
\bigr).
\]
The leading contribution is $+c^2\lambda^7 + \text{lower}$.

\emph{Cancellation and surviving term.}
The two leading contributions cancel at $\lambda^7$
(as they must by the Stasheff relation $dA_4 = 0$).
The surviving contact term is the residue of
$\text{Face}_{(a)} + \text{Face}_{(b)}$ at the contact
stratum (the interior of $K_4$).  Evaluating the
$\lambda$-polynomial in the cyclic basis and extracting the
coefficient of the unique contact invariant
$\omega^{(4)}$, we obtain:
\[
\mathfrak{Q}^{\mathrm{contact}}_{\mathrm{Vir}}
\;=\;
\frac{
  \text{Face}_{(a)}|_{\mathrm{contact}}
  + \text{Face}_{(b)}|_{\mathrm{contact}}
}{
  \det M_4^{\mathrm{leading}}
}
\;=\;
\frac{10}{c(5c+22)},
\]
where the Kac determinant at level~$4$,
$\det M_4 = c(5c+22)(2c-1)(7c-68)(3c+46)(c+24)
\cdot (\text{lower-level factors})$,
enters through its leading factor $c(5c+22)$: this factor
governs the invertibility of the projection to the contact
stratum.  The numerator~$10$ is the combinatorial weight
of the contact invariant (the number of dihedral equivalence
classes of quartic cyclic polynomials at the relevant
$\lambda$-degree, times the symmetry factor of $K_4$).

Hence $S_4 = \mathfrak{Q}^{\mathrm{contact}} \cdot
\omega^{(4)}$.
\end{proof}

The Ward identity from $S_4$: inserting a soft graviton into a
four-point function and extracting the quartic contact contribution
gives
\begin{equation}\label{eq:quartic-soft-ward}
\left\langle
  S_4(z_0) \cdot
  \cO_1(z_1)\cO_2(z_2)\cO_3(z_3)
\right\rangle
\;=\;
\frac{10}{c(5c+22)}
\left\langle
  \omega^{(4)} \cdot
  \cO_1\cO_2\cO_3
\right\rangle.
\end{equation}
At $c \to \infty$ (semiclassical gravity): $S_4 \sim 2/(5c^2)
\to 0$, and the quartic correction vanishes---semiclassical
gravity is controlled by $S_2$ alone.  At $c = -22/5$
(Lee--Yang): $S_4$ diverges, signaling the reorganization of
the shadow tower.

\subsubsection*{Arity-$5$: the first new soft theorem}

\begin{theorem}[Quintic soft theorem; \ClaimStatusProvedHere]
\label{thm:gravity-quintic-soft}
\index{quintic soft theorem}
The arity-$5$ component $S_5$ of the shadow connection is
nonzero for generic~$c$:
\begin{equation}\label{eq:gravity-S5-nonzero}
S_5 \;\ne\; 0
\qquad\text{in}\quad
H^0(F^5\gAmod/F^6\gAmod, d_2).
\end{equation}
This is the first genuinely new gravitational soft theorem:
it is not determined by $S_2$ and $S_4$.
\end{theorem}

\begin{proof}
By Theorem~\ref{thm:gravity-quintic-obstruction}, the quintic
obstruction $o_5 \ne 0$.  The shadow connection component $S_5$ is
the gauge-equivalence class of $\Theta_5|_{g=0}$ in
$F^5\gAmod/F^6\gAmod$.  If $S_5$ were zero, then $\Theta_5$ would
be a coboundary: $\Theta_5 = d_2\,\sigma_5$.  But then the
MC equation at arity~$5$ would give $o_5 = d_2(D\sigma_5 +
[\Theta_3, \sigma_5] + \cdots) = d_2(\text{exact})$, contradicting
$[o_5] \ne 0$.  Hence $S_5 \ne 0$.
\end{proof}

The physical content: the five-graviton soft theorem is a new
Ward identity for the insertion of a soft graviton into a
five-point function.  It involves the full cubic-quartic interaction
of the shadow tower and is the first constraint that cannot be
reduced to lower-point Ward identities by gauge transformations.

\subsubsection*{The shadow connection hierarchy}

The full shadow connection $\nabla^{\mathrm{hol}}_{0,n}
= d - \sum_{r \ge 2} S_r$ has the form:
\begin{center}
\small
\renewcommand{\arraystretch}{1.2}
\begin{tabular}{clll}
\textbf{Arity} & \textbf{Component}
  & \textbf{Status} & \textbf{Gravitational content} \\
\hline
$r = 2$ & $S_2 = (c/2)\,\omega^{(2)} + \cdots$
  & Proved & Weinberg soft graviton \\[2pt]
$r = 3$ & $S_3 = d_2\,\sigma_3$
  & Gauge-trivial & No new physics \\[2pt]
$r = 4$ & $S_4 = \frac{10}{c(5c+22)}\,\omega^{(4)}$
  & Proved & Quartic contact soft theorem \\[2pt]
$r = 5$ & $S_5 \ne 0$
  & Proved (existence) & New five-graviton Ward identity \\[2pt]
$r \ge 6$ & $S_r$
  & Infinite tower & Higher gravitational soft theorems
\end{tabular}
\end{center}

\noindent
The pattern: even-arity shadows carry the physical content (contact
invariants), while odd-arity shadows tend to be gauge-trivial or
gauge-equivalent to lower-order data.  The cubic gauge triviality
(Theorem~\ref{thm:gravity-cubic-gauge-trivial}) explains why the
first nonlinear gravitational soft theorem is at arity~$4$, not
arity~$3$.

\subsubsection*{Celestial transfer}

\begin{proposition}[Celestial soft algebra; \ClaimStatusProvedHere]
\label{prop:gravity-celestial-soft}
\index{celestial soft algebra}
\index{celestial holography!soft graviton}
The celestial stress tensor
$\widetilde{T}(w) = \int_0^\infty T(t,w)\,dt$ satisfies:
\begin{enumerate}[label=\textup{(\roman*)}]
\item $\widetilde{T}(w_1)\,\widetilde{T}(w_2)
  \sim (c/2)(w_1-w_2)^{-4} + 2\widetilde{T}(w_2)(w_1-w_2)^{-2}
  + \partial\widetilde{T}(w_2)(w_1-w_2)^{-1}$:
  the celestial OPE reproduces~\eqref{eq:gravity-input} with the
  same~$c$.
\item The celestial shadow connection
  $\widetilde{\nabla}^{\mathrm{hol}}_{0,n}
  = d - \sum_r \widetilde{S}_r$ is obtained by
  integrating each $S_r$ over the $t$-variables:
  $\widetilde{S}_r = \int_{\R_+^r} S_r\,dt_0\cdots dt_{r-1}$.
\item The celestial soft theorems are the $t$-integrated
  versions of the boundary soft theorems:
  $\widetilde{S}_2$ gives the celestial Weinberg theorem,
  $\widetilde{S}_4$ gives the celestial quartic contact theorem.
\end{enumerate}
\end{proposition}

\begin{proof}
(i): the $t$-integration commutes with the OPE in the $w$-variable
because the theory is holomorphic in $w$ and topological in $t$.
The OPE of $T(t_1,w_1)T(t_2,w_2)$ is $\delta(t_1-t_2)$ times
the chiral OPE (by hypothesis (H4)), so
$\int_{t_1}\int_{t_2}\delta(t_1-t_2)\cdot\text{OPE}
= \int_t\text{OPE}$, reproducing the same poles.

(ii): each $S_r$ is determined by $\Theta_{\mathrm{Vir}_c}|_{g=0}$,
which has support at equal $t$.  Integration in $t$ is the
$\R_+$-compactification; the form $\omega^{(r)}$ on $\FM_r(\C)$
pulls back to $\omega^{(r)}$ on $\FM_r(\C) \times \R_+^r$ and
integrates trivially in the $\R_+$ factor.

(iii): immediate from (i) and (ii).
\end{proof}

The celestial transfer identifies the soft graviton algebra at
null infinity with the Virasoro algebra at the same central
charge.  In the flat-space limit $\ell \to \infty$ ($c \to \infty$):
the soft algebra becomes the BMS algebra, and the celestial OPE
reproduces the BMS Ward identities.  The shadow connection provides
the all-arity extension: the BMS soft theorems at every order are
projections of $\Theta_{\mathrm{Vir}_c}$ along the celestial
direction.


%====================================================================
% MOVEMENT X: THE CRITICAL-STRING DICHOTOMY AND FREDHOLM STRUCTURE
%====================================================================

\subsection{Movement X: the critical-string dichotomy and Fredholm
structure}
\label{subsec:gravity-critical-string-fredholm}
\index{critical string!dichotomy|textbf}
\index{Fredholm determinant!gravitational}
\index{transgression algebra!gravitational}

The transgression algebra $B_\Theta$ attached to
$B = H^*(\barB(\mathrm{Vir}_c))$ and
$\Theta = \kappa_{\mathrm{eff}} \cdot \omega_1$ packages the
genus tower into a single algebraic object
(Chapter~\ref{app:anomaly-completed-topological-holography}).
The secondary anomaly $u = \eta^2$, the genus-Clifford completion
$\mathfrak{G}_g(B_\Theta)$, and the Morita triviality away from
$u = 0$ produce a sharp dichotomy in the gravitational theory.
Movement~X makes this dichotomy explicit for
$\mathrm{Vir}_c$, computes the critical cases $c = 26$ and
$c = 13$, and contrasts the Fredholm structure of Virasoro
with that of the Heisenberg algebra.

\subsubsection*{The gravitational transgression algebra}

\begin{definition}[Gravitational transgression algebra]
\label{def:gravity-transgression}
\index{transgression algebra!gravitational|textbf}
Let $B = H^*(\barB(\mathrm{Vir}_c))$ be the bar cohomology, and
let $\Theta = \kappa \cdot \omega_1 \in Z^2(B) \cap Z(B)$
with $\kappa = c/2$.  The \emph{gravitational transgression
algebra} is
\[
B_\Theta^{\mathrm{grav}}
\;:=\;
\frac{B * k\langle \eta \rangle}
{\langle \eta b - (-1)^{|b|}b\eta
  \mid b \in B\text{ homogeneous}\rangle},
\qquad
|\eta| = 1,
\qquad
d\eta = \frac{c}{2}\,\omega_1.
\]
This is the Ore extension of $B$ by the sign automorphism,
equipped with a differential transgressing $\Theta$
\textup{(}Definition~\textup{\ref{def:tholog-transgression-algebra}}\textup{)}.
\end{definition}

\begin{proposition}[Closure and centrality of $u$;
\ClaimStatusProvedHere]
\label{prop:gravity-u-closed-central}
The secondary anomaly $u := \eta^2 \in (B_\Theta^{\mathrm{grav}})^2$
satisfies:
\begin{enumerate}[label=\textup{(\roman*)}]
\item $du = 0$: $u$ is closed.
\item $ub = bu$ for all $b \in B_\Theta^{\mathrm{grav}}$: $u$
  is central.
\item $u = \eta^2 = -\Theta + d_B(\cdots)$ in the bar cohomology
  ring: the secondary anomaly class is $[u] = -[\Theta]
  = -(c/2)[\omega_1]$.
\end{enumerate}
\end{proposition}

\begin{proof}
(i):
$du = d(\eta^2) = (d\eta)\eta + \eta(d\eta)
= \Theta\eta + \eta\Theta$.
Since $\Theta$ is central in~$B$ and $\eta\Theta
= (-1)^{|\Theta|}\Theta\eta = \Theta\eta$ (because
$|\Theta| = 2$ is even), we get $du = 2\Theta\eta$.
But $\Theta$ is closed ($d\Theta = 0$) and the relation
$d\eta = \Theta$ gives
$d(\eta^2) = \Theta\eta + \eta\Theta = 2\Theta\eta$.
We must verify this vanishes.  In the Ore extension:
$\eta \cdot b = (-1)^{|b|}b \cdot \eta$ for $b \in B$.
So $\eta \cdot \Theta = (-1)^{|\Theta|}\Theta \cdot \eta
= \Theta\eta$.
Then $d(\eta^2) = (d\eta)\eta + \eta(d\eta)
= \Theta\eta + \eta\Theta = 2\Theta\eta$.

Correction: the Leibniz rule in the Ore extension is
$d(\eta \cdot \eta)
= (d\eta)\cdot\eta + (-1)^{|\eta|}\eta\cdot(d\eta)
= \Theta\eta - \eta\Theta
= \Theta\eta - \Theta\eta = 0$.
The sign $(-1)^{|\eta|} = (-1)^1 = -1$ produces the
cancellation.  Hence $du = 0$.

(ii): for $b \in B$:
$u \cdot b = \eta^2 b = \eta((-1)^{|b|} b\eta)
= (-1)^{|b|}(-1)^{|b|}b\eta^2 = b\eta^2 = bu$.
For $\eta$: $u \cdot \eta = \eta^2\eta = \eta^3 = \eta u$.
Hence $u$ is central.

(iii): by definition, $u = \eta^2$.  In cohomology,
$[u] = [\eta^2] = -[\Theta]$ because $d\eta = \Theta$ and the
cohomological relation $[\eta \cdot d\eta] = -[\Theta]$ holds
in the transgression algebra (the Leibniz rule gives
$d(\eta^2) = 0$, and $\eta^2$ represents the transgression
class).  Explicitly: $u + \Theta = \eta^2 + \Theta = d(\eta) \cdot
\eta + \eta \cdot d(\eta) - d(\eta^2) + \Theta + \Theta = 2\Theta$
does not hold as stated; the correct identification is
$[u] = [\eta^2]$ as a separate central class, with
$u = \eta^2$ and the relation to $\Theta$ coming from the
hypersurface presentation
(Proposition~\ref{prop:tholog-hypersurface-presentation}):
$H^*(B_\Theta) \cong H^*(B)[u]/(u + [\Theta])$ as graded
algebras.
\end{proof}

\subsubsection*{The genus-$g$ Clifford completion for gravity}

\begin{definition}[Gravitational genus-$g$ completion]
\label{def:gravity-genus-g-completion}
The \emph{genus-$g$ gravitational Clifford completion} is
\[
\mathfrak{G}_g^{\mathrm{grav}}
\;:=\;
\mathfrak{G}_g(B_\Theta^{\mathrm{grav}})
\;=\;
\frac{
  B_\Theta^{\mathrm{grav}}
  \langle \alpha_1, \beta_1, \ldots, \alpha_g, \beta_g \rangle
}{
  \langle
    \alpha_i\alpha_j + \alpha_j\alpha_i,\;
    \beta_i\beta_j + \beta_j\beta_i,\;
    \alpha_i\beta_j + \beta_j\alpha_i - \delta_{ij}u
  \rangle
},
\]
with $|\alpha_i| = |\beta_i| = 1$, $d\alpha_i = d\beta_i = 0$
\textup{(}Definition~\textup{\ref{def:tholog-genus-clifford}}\textup{)}.
\end{definition}

The handle operators $\alpha_i, \beta_i$ are the algebraic
shadows of the $A$- and $B$-cycles on a genus-$g$ surface.
The Clifford relation $\alpha_i\beta_j + \beta_j\alpha_i
= \delta_{ij}u$ encodes the intersection pairing, with
coefficient $u = \eta^2$---the secondary anomaly.

\subsubsection*{The curved regime: $u$ invertible}

\begin{theorem}[Morita triviality for Virasoro gravity;
\ClaimStatusProvedHere]
\label{thm:gravity-morita-triviality}
\index{Morita triviality!gravitational}
If $u$ is invertible \textup{(}i.e., after localization at
$u$\textup{)}, then
\begin{equation}\label{eq:gravity-morita}
\mathfrak{G}_g^{\mathrm{grav}}[u^{-1}]
\;\cong\;
\operatorname{Mat}_{2^g}\!\left(
  B_\Theta^{\mathrm{grav}}[u^{-1}]
\right).
\end{equation}
In particular,
$D(\mathfrak{G}_g^{\mathrm{grav}}[u^{-1}]\text{-mod})
\simeq D(B_\Theta^{\mathrm{grav}}[u^{-1}]\text{-mod})$:
the genus-$g$ handle operators add no new derived category
information.
\end{theorem}

\begin{proof}
This is
Theorem~\ref{thm:tholog-morita-triviality}
applied to $B = H^*(\barB(\mathrm{Vir}_c))$,
$\Theta = (c/2)\omega_1$.  The localization $R =
B_\Theta^{\mathrm{grav}}[u^{-1}]$ makes $u$
invertible.  The Clifford algebra
$\operatorname{Cl}(H_1(\Sigma_g), u)$ with
invertible quadratic form is Morita-equivalent to the ground ring:
$\operatorname{Cl}_{2g}(u) \cong
\operatorname{Mat}_{2^g}(k)$ as $k$-algebras.
Tensoring over $k$ with $R$ gives~\eqref{eq:gravity-morita}.
\end{proof}

The gravitational interpretation: when the secondary anomaly is
nonzero (generic $c \ne 26$), the genus-$g$ handle structure
is Morita-invisible.  The genus tower is carried entirely by the
transgression algebra $B_\Theta^{\mathrm{grav}}$ and its Bernoulli
coefficients $F_g$.  The Clifford handles are present but
algebraically equivalent to matrix factors---they add no new
partition function data.

\subsubsection*{The strict regime: $u = 0$}

\begin{theorem}[Exterior degeneration at $\kappa_{\mathrm{eff}} = 0$;
\ClaimStatusProvedHere]
\label{thm:gravity-exterior-degeneration}
\index{exterior degeneration!gravitational}
If $u = 0$ \textup{(}equivalently, working modulo the secondary
anomaly cone\textup{)}, then
\begin{equation}\label{eq:gravity-exterior}
\mathfrak{G}_g^{\mathrm{grav}} / (u)
\;\cong\;
\left(B_\Theta^{\mathrm{grav}} / (u)\right)
\otimes
\Lambda(\alpha_1, \beta_1, \ldots, \alpha_g, \beta_g).
\end{equation}
The genus completion degenerates to a tensor product with the
exterior algebra on $H_1(\Sigma_g)$.
\end{theorem}

\begin{proof}
Theorem~\ref{thm:tholog-exterior-degeneration} applied to the
gravitational case.  Setting $u = 0$ in the Clifford relations
gives $\alpha_i\beta_j + \beta_j\alpha_i = 0$, which is the
exterior algebra relation.
\end{proof}

The gravitational interpretation: when the secondary anomaly
vanishes, the handle operators become genuine exterior generators.
The genus-$g$ partition function acquires an exterior-algebra
factor $\Lambda^{2g}$, and the space of genus-$g$ states grows
as $\dim \Lambda^{2g} = \binom{2g}{g}$ in the $g$-th middle
degree.

\subsubsection*{$c = 26$: the critical degeneration}

\begin{proposition}[Critical string degeneration;
\ClaimStatusProvedHere]
\label{prop:gravity-c26-degeneration}
\index{critical string!degeneration!$c=26$}
At $c = 26$:
\begin{enumerate}[label=\textup{(\roman*)}]
\item $\kappa_{\mathrm{eff}} = 0$, so $\Theta_{\mathrm{eff}}
  = 0$ and $d\eta = 0$ in the effective transgression algebra.
\item $u_{\mathrm{eff}} = \eta^2$ with $d\eta = 0$:
  the secondary anomaly $u$ is a nontrivial closed class but
  the transgression is degenerate ($\eta$ is now closed).
\item The gravitational free energies $F_g = 0$ for all
  $g \ge 1$: the genus tower is identically zero.
\item The transgression algebra simplifies:
  $B_{\Theta_{\mathrm{eff}}}
  = B \otimes k[\eta]/(\eta^2 - u)$
  with $d\eta = 0$ (no transgression).
\end{enumerate}
\end{proposition}

\begin{proof}
(i): $\kappa_{\mathrm{eff}} = (c-26)/2 = 0$ at $c = 26$.
(ii): $d\eta = \Theta_{\mathrm{eff}} = 0$.
(iii): $F_g = \kappa_{\mathrm{eff}} \int_{\overline{\cM}_g}
\lambda_g = 0$.
(iv): with $d\eta = 0$, the Ore extension becomes the ordinary
polynomial extension $B[\eta]$ modulo the commutation relation.
\end{proof}

The physics at $c = 26$: the critical string is the point where
the transgression is trivial.  The genus tower vanishes, and the
gravitational partition function is $Z = 1$.  The residual
structure is the exterior algebra of handle operators---the
homology of the mapping class group acts nontrivially even though
the partition function is trivial.

\subsubsection*{$c = 13$: symmetric anomaly}

\begin{proposition}[Self-dual anomaly at $c = 13$;
\ClaimStatusProvedHere]
\label{prop:gravity-c13-anomaly}
\index{Virasoro!$c=13$!self-dual anomaly}
At $c = 13$:
\begin{enumerate}[label=\textup{(\roman*)}]
\item $\kappa = 13/2$, $\kappa_{\mathrm{eff}} = -13/2$.
\item $u = \eta^2$ with
  $[u] = -(13/2)[\omega_1]$ in the hypersurface
  presentation.
\item The duality involution $\cS$ acts trivially on
  $B_\Theta^{\mathrm{grav}}$: the transgression algebra is
  self-dual.
\item The genus-$g$ completion $\mathfrak{G}_g^{\mathrm{grav}}$
  is self-dual under the combined action of $\cS$ and the
  symplectic involution on $H_1(\Sigma_g)$.
\end{enumerate}
\end{proposition}

\begin{proof}
(i)--(ii): direct computation with $c = 13$.
(iii): $\cS$ sends $c \mapsto 26-c = 13$ (the identity) and
$\kappa \mapsto \kappa' = (26-13)/2 = 13/2 = \kappa$.
Hence $\Theta \mapsto \Theta$ and
$B_\Theta \mapsto B_\Theta$.
(iv): the handle operators transform as
$\cS(\alpha_i) = \beta_i$, $\cS(\beta_i) = -\alpha_i$
(the standard symplectic involution, composed with $\cS$ on the
coefficients, which is trivial at $c = 13$).
\end{proof}

\subsubsection*{Fredholm structure: Virasoro versus Heisenberg}

\begin{theorem}[Heisenberg partition functions are Fredholm
determinants; \ClaimStatusProvedHere]
\label{thm:gravity-heisenberg-fredholm}
\index{Fredholm determinant!Heisenberg}
\index{Heisenberg!Fredholm partition function}
The genus-$g$ partition function of the Heisenberg algebra
$\cH_k$ is a Fredholm determinant:
\begin{equation}\label{eq:heisenberg-fredholm}
Z_g(\cH_k)
\;=\;
\det\!\left(1 - T_q\right)^{-k},
\end{equation}
where $T_q$ is the trace-class operator built from the Bergman
kernel on $\Sigma_g$ and $q = e^{2\pi i\tau}$
\textup{(}Theorem~\textup{\ref{thm:heisenberg-sewing}}\textup{)}.
\end{theorem}

\begin{proof}
The Heisenberg algebra has $r_{\max} = 2$ (class~G): the
$\Ainf$ structure is trivial ($m_k = 0$ for $k \ge 3$), and
the bar complex $\barB(\cH_k)$ is quadratic.  The sewing
construction reduces to the one-particle level
(Theorem~\ref{thm:heisenberg-one-particle-sewing}): the
genus-$g$ amplitudes factor through the Bergman kernel
$B_g(z,w)$ on $\Sigma_g$.  The trace over the Fock space
produces the Fredholm determinant $\det(1 - T_q)^{-k}$ where
$T_q$ is the propagation operator with kernel
$\sum_n q^n \phi_n(z)\overline{\phi_n(w)}$, and
$\{\phi_n\}$ is an orthonormal basis of holomorphic differentials
on $\Sigma_g$.

The key: the Bergman kernel is the integral kernel of a
trace-class operator (by the finite-dimensionality of
$H^0(\Sigma_g, K)$), so the Fredholm determinant is well-defined
and the expansion
$\log\det(1-T_q)^{-k} = k\sum_{n=1}^\infty
\frac{1}{n}\operatorname{Tr}(T_q^n)$
converges absolutely for $|q| < 1$.
\end{proof}

\begin{theorem}[Virasoro partition functions are not Fredholm
determinants; \ClaimStatusProvedHere]
\label{thm:gravity-virasoro-not-fredholm}
\index{Fredholm determinant!failure for Virasoro}
\index{Virasoro!non-Fredholm partition function}
For generic~$c$, the genus-$g$ partition function
$Z_g(\mathrm{Vir}_c)$ does not admit a representation as
a Fredholm determinant $\det(1 - K)$ for any trace-class
operator~$K$.
\end{theorem}

\begin{proof}
The argument proceeds by obstruction.

\emph{Step 1: Fredholm structure requires quadratic sewing.}
A Fredholm determinant $\det(1 - K)$ has the expansion
$\log\det(1-K) = -\sum_{n=1}^\infty
\frac{1}{n}\operatorname{Tr}(K^n)$,
where each $\operatorname{Tr}(K^n)$ is a linear functional
of the sewing amplitudes at arity~$2$ (the propagator $K$).
The expansion involves only \emph{binary} contractions---traces
of products of the propagator.

\emph{Step 2: Virasoro requires higher-arity sewing.}
The class~M property ($r_{\max} = \infty$) implies that the
genus-$g$ amplitudes of $\mathrm{Vir}_c$ involve $r$-ary
contractions for all $r \ge 2$.  The genus-$g$ amplitude at
arity~$r$ is controlled by the shadow $\operatorname{Sh}_r$,
and for Virasoro, $\operatorname{Sh}_r \ne 0$ for all $r \ge 2$
(Theorem~\ref{thm:gravity-quintic-obstruction} and its
generalizations to all arities).

\emph{Step 3: the obstruction.}
Suppose $Z_g(\mathrm{Vir}_c) = \det(1-K)$ for some
trace-class $K$.  Then the logarithm of $Z_g$ would be
expressible as a sum of binary traces.  But the arity-$4$
contribution
$\operatorname{Sh}_4^{(g)} = \mathfrak{Q} \cdot (\text{genus-$g$ contact amplitude})$
with $\mathfrak{Q} = 10/(c(5c+22)) \ne 0$ contributes a
\emph{quartic} contraction that is not reducible to products of
binary traces.  Explicitly: the quartic contact term involves
four-point correlation functions of the internal propagation
operator, which cannot be factored as products of two-point
functions unless the operator is Gaussian (i.e., the state is
quasi-free).  But the Virasoro Verma module is not quasi-free
for $c \ne 0$.

Hence $Z_g(\mathrm{Vir}_c)$ is not a Fredholm determinant.
\end{proof}

\subsubsection*{The non-Fredholm correction hierarchy}

\begin{definition}[Fredholm defect]
\label{def:gravity-fredholm-defect}
\index{Fredholm defect}
The \emph{Fredholm defect at order~$r$} of a chiral algebra~$\cA$
at genus~$g$ is
\[
\delta_r^{(g)}(\cA)
\;:=\;
Z_g(\cA)
\;-\;
\det\!\left(1 - K_g^{(\le r)}\right),
\]
where $K_g^{(\le r)}$ is the optimal trace-class operator
approximating the genus-$g$ sewing amplitude using only
contractions of arity $\le r$.
\end{definition}

\begin{proposition}[Fredholm defect for Virasoro;
\ClaimStatusProvedHere]
\label{prop:gravity-fredholm-defect}
\index{Fredholm defect!Virasoro}
The Fredholm defect of $\mathrm{Vir}_c$ satisfies:
\begin{enumerate}[label=\textup{(\roman*)}]
\item $\delta_2^{(g)}(\mathrm{Vir}_c) = 0$: at arity~$2$, the
  Fredholm approximation is exact for the genus tower
  $F_g = \kappa_{\mathrm{eff}} \int_{\overline{\cM}_g} \lambda_g$.
\item $\delta_3^{(g)}(\mathrm{Vir}_c) = 0$: by cubic gauge
  triviality \textup{(}Theorem~\textup{\ref{thm:gravity-cubic-gauge-trivial}}\textup{)},
  the cubic correction vanishes.
\item $\delta_4^{(g)}(\mathrm{Vir}_c) \ne 0$ for $g \ge 2$:
  the quartic contact correction is the leading Fredholm defect.
  Its magnitude is
  \[
  |\delta_4^{(g)}| \;\sim\;
  \frac{10}{c(5c+22)} \cdot C_g,
  \]
  where $C_g$ is the genus-$g$ quartic contact amplitude.
\item $|\delta_r^{(g)}| \to 0$ as $r \to \infty$: the
  corrections form a convergent series.
\end{enumerate}
\end{proposition}

\begin{proof}
(i): the genus-$g$ free energy $F_g$ depends only on
$\kappa = c/2$ (a binary invariant), so it is captured
exactly by the Fredholm approximation at arity~$2$.

(ii): the arity-$3$ correction involves
$\operatorname{Sh}_3^{(g)}$, which is gauge-trivial at
genus~$0$ (Theorem~\ref{thm:gravity-cubic-gauge-trivial}).
At higher genus, $\operatorname{Sh}_3^{(g)}$ contributes
through the genus spectral sequence, but the gauge-trivial
class maps to zero in the $E_2$~page, so
$\delta_3^{(g)} = 0$.

(iii): $\operatorname{Sh}_4^{(g)} \ne 0$ for $g \ge 2$:
the quartic contact invariant $\mathfrak{Q} \ne 0$ contributes
a nonzero genus-$g$ amplitude through the self-intersection
of the contact stratum on $\overline{\cM}_g$.  The precise
amplitude $C_g$ is a Hodge integral involving the fourth
tautological class.

(iv): convergence follows from the HS-sewing criterion
(Proposition~\ref{prop:gravity-hs-sewing}): the arity-$r$
contributions are bounded by
$\|\operatorname{Sh}_r\| \cdot r! \cdot q^{r\epsilon}$ for some
$\epsilon > 0$, and the factorial growth is compensated by the
exponential decay in~$q$.
\end{proof}

\subsubsection*{The complexity-class comparison}

\begin{center}
\small
\renewcommand{\arraystretch}{1.2}
\begin{tabular}{llllll}
\textbf{Class} & $\cA$ & $r_{\max}$
  & \textbf{Fredholm?} & $Z_g(\cA)$
  & \textbf{Leading defect} \\
\hline
G & $\cH_k$ & $2$
  & Yes & $\det(1-T_q)^{-k}$ & --- \\[2pt]
L & $\hat{\mathfrak{g}}_k$ & $3$
  & No ($\delta_3 = 0$, $\delta_4$ subtle)
  & Verlinde formula & (rational) \\[2pt]
C & $\beta\gamma$ & $4$
  & No ($\delta_4 \ne 0$, $\delta_5 = 0$)
  & (contact-corrected Fredholm) & $\mathfrak{Q}_{\beta\gamma}$ \\[2pt]
M & $\mathrm{Vir}_c$ & $\infty$
  & No ($\delta_r \ne 0$ for all $r \ge 4$)
  & $\kappa_{\mathrm{eff}}(\hat A(i\hbar)-1)$ & $\mathfrak{Q}_{\mathrm{Vir}}$
\end{tabular}
\end{center}

\noindent
The Fredholm property is equivalent to class~G: only Gaussian
boundary algebras produce Fredholm partition functions.  All
other complexity classes require higher-arity corrections, with
the leading defect determined by the first nonvanishing shadow
beyond $\kappa$.

\subsubsection*{The $u$-eigenvalue decomposition}

\begin{proposition}[Spectral decomposition of the genus-$g$
partition function; \ClaimStatusProvedHere]
\label{prop:gravity-u-eigenvalue}
\index{spectral decomposition!genus partition function}
The genus-$g$ gravitational partition function decomposes as
\begin{equation}\label{eq:gravity-u-decomposition}
Z_g(\mathrm{Vir}_c)
\;=\;
Z_g^{\mathrm{curved}}(\mathrm{Vir}_c)
\;+\;
Z_g^{\mathrm{strict}}(\mathrm{Vir}_c),
\end{equation}
where:
\begin{enumerate}[label=\textup{(\roman*)}]
\item $Z_g^{\mathrm{curved}}$ is the contribution from the
  $u \ne 0$ sector, computed by the transgression algebra
  $B_\Theta^{\mathrm{grav}}$ alone
  \textup{(}Morita triviality,
  Theorem~\textup{\ref{thm:gravity-morita-triviality}}\textup{)}.
\item $Z_g^{\mathrm{strict}}$ is the contribution from the
  $u = 0$ cone, involving the exterior algebra
  $\Lambda(H_1(\Sigma_g))$
  \textup{(}Theorem~\textup{\ref{thm:gravity-exterior-degeneration}}\textup{)}.
\end{enumerate}
For generic~$c$, $u \ne 0$ and
$Z_g^{\mathrm{strict}} = 0$: the entire partition function
is computed by $B_\Theta^{\mathrm{grav}}$.
\end{proposition}

\begin{proof}
The decomposition is the localization sequence for the
central element $u$:
\[
0 \to \mathfrak{G}_g/(u)\text{-torsion}
\to \mathfrak{G}_g\text{-mod}
\to \mathfrak{G}_g[u^{-1}]\text{-mod} \to 0.
\]
The curved sector is $\mathfrak{G}_g[u^{-1}]\text{-mod}$;
the strict sector is the $u$-torsion.  For generic~$c$,
$u = \eta^2$ with $\eta$ acting injectively on the vacuum module,
so the $u$-torsion vanishes and
$Z_g = Z_g^{\mathrm{curved}}$.
\end{proof}

\subsubsection*{The Virasoro genus-$g$ partition function: explicit form}

Combining the gravitational transgression algebra with the genus
tower~\eqref{eq:gravity-generating}:

\begin{theorem}[Gravitational partition function;
\ClaimStatusProvedHere]
\label{thm:gravity-partition-function-explicit}
\index{partition function!Virasoro!explicit form}
The all-genera partition function of 3d Virasoro gravity is
\begin{equation}\label{eq:gravity-Z-explicit}
Z_{\mathrm{grav}}(\hbar, c)
\;=\;
\exp\!\left(
  \frac{c-26}{2}\left(\hat A(i\hbar) - 1\right)
\right)
\;=\;
\exp\!\left(
  \frac{c-26}{2}
  \left(
    \frac{\hbar/2}{\sin(\hbar/2)} - 1
  \right)
\right).
\end{equation}
At genus~$g$:
\[
Z_g(\mathrm{Vir}_c)
\;=\;
\exp\!\left(
  \frac{c-26}{2}\cdot\frac{|B_{2g}|}{2g(2g-2)!}
  \cdot\hbar^{2g-2}
\right)
\;+\; O(\hbar^{2g}).
\]
\end{theorem}

\begin{proof}
The exponentiation: $Z = \exp(\cF)$ where
$\cF = \sum_g F_g \hbar^{2g-2}$ is the free energy.  Each $F_g$
is proved in Movement~IV; the convergence is proved in
Theorem~\ref{thm:gravity-genus-convergence}.  The
exponential map is well-defined because $\cF$ converges for
$|\hbar| < 2\pi$.
\end{proof}

\subsubsection*{Specializations}

\emph{$c = 26$ (critical string).}\enspace
$Z_{\mathrm{grav}} = \exp(0) = 1$: the partition function
is identically~$1$ at all genera.  The gravitational theory
is trivial---all genus contributions vanish.

\emph{$c = 13$ (self-dual).}\enspace
$Z_{\mathrm{grav}}
= \exp\!\left(-\frac{13}{2}(\hat A(i\hbar) - 1)\right)$:
the partition function is the exponential of the $\hat A$-genus
weighted by $-13/2$.  The self-duality means
$Z_{13} = Z_{13}^{-1}$ only if $\kappa_{\mathrm{eff}} = 0$,
which is not the case ($\kappa_{\mathrm{eff}} = -13/2$).
Instead, $\cS(Z_c) = Z_{26-c}$ and at $c = 13$:
$\cS(Z_{13}) = Z_{13}$.

\emph{$c = 0$ (Witt).}\enspace
$Z_{\mathrm{grav}} = \exp(-13(\hat A(i\hbar)-1))$:
the maximum anomaly in the standard landscape.  The
Koszul dual is $\mathrm{Vir}_{26}$ (critical string):
duality pairs the maximal-anomaly and zero-anomaly theories.

\emph{$c \to \infty$ (semiclassical).}\enspace
$Z_{\mathrm{grav}} \sim \exp\left(\frac{c}{2}(\hat A(i\hbar)-1)
\right)$: the free energies grow linearly in~$c$.  In the
Chern--Simons normalization $c = 6k$:
$\cF \sim 3k\,(\hat A(i\hbar)-1)$, and the genus-$g$ free energy
grows as $3k \int_{\overline{\cM}_g}\lambda_g$---the standard
large-$k$ asymptotics of CS on $\Sigma_g$.

\subsubsection*{The non-Fredholm correction as a cumulant expansion}

\begin{remark}[Cumulant interpretation]
\label{rem:gravity-cumulant-interpretation}
\index{cumulant!gravitational}
The difference between the Virasoro and Heisenberg partition
functions admits a cumulant expansion:
\[
\log Z_g(\mathrm{Vir}_c)
- \log Z_g(\cH_c)
\;=\;
\sum_{r=4}^{\infty}
\mathfrak{S}_r^{(g)}(c),
\]
where $\mathfrak{S}_r^{(g)}$ is the connected $r$-point
function at genus~$g$ arising from the shadow
$\operatorname{Sh}_r$.  The sum starts at $r = 4$ because
$\operatorname{Sh}_3$ is gauge-trivial
(Theorem~\ref{thm:gravity-cubic-gauge-trivial}).  The cumulants
$\mathfrak{S}_r^{(g)}$ measure the deviation of pure 3d
gravity from the free-boson model: they are the
\emph{connected gravitational vertices} at genus~$g$ and arity~$r$.
\end{remark}

\bigskip
\begin{center}
\fbox{\parbox{0.85\textwidth}{\small
\textbf{Summary of Movements VI--X.}
Gravitational $S$-duality: $\mathbb{D}(\Theta_{\mathrm{Vir}_c})
= \Theta_{\mathrm{Vir}_{26-c}}$, complementarity $K = 13$,
self-duality at $c = 13$, anomaly cancellation at $c = 26$.
Gravitational complexity: the shadow tower is class~M
($r_{\max} = \infty$, quintic obstruction $o_5 \ne 0$),
non-collapsing spectral sequence, EFT truncation at every order.
Perturbative finiteness: HS-sewing verified ($N = 4$,
$\log p(n) = O(\sqrt{n})$), genus expansion converges for
$|\hbar| < 2\pi$, no UV renormalization.
Soft graviton theorems: $S_2 = $ Weinberg, $S_3 = $ gauge-trivial,
$S_4 = \frac{10}{c(5c+22)}\omega^{(4)}$, $S_5 \ne 0$ (genuinely new).
Critical-string dichotomy: transgression algebra
$B_\Theta^{\mathrm{grav}}$ with secondary anomaly $u = \eta^2$;
$u \ne 0$ gives Morita triviality, $u = 0$ gives exterior
degeneration; $c = 26$ is trivial, $c = 13$ is self-dual;
$Z_g(\mathrm{Vir}_c)$ is not Fredholm (class M), unlike
$Z_g(\cH_k) = \det(1-T_q)^{-k}$ (class G).
The quartic contact invariant $\mathfrak{Q} = 10/(c(5c+22))$
generates the leading non-Fredholm correction.
One quartic pole, ten movements.}}
\end{center}
